\documentclass[12pt,a4paper]{report}

% --- Packages ----------------------------------------------------------------
\usepackage[utf8]{inputenc}
\usepackage[T1]{fontenc}
\usepackage[italian]{babel}
\usepackage{geometry}
\usepackage{graphicx}
\usepackage{amsmath}
\usepackage{hyperref}
\usepackage{booktabs}
\usepackage{longtable}
\usepackage{array}
\usepackage{xcolor}
\usepackage{listings}
\usepackage{fancyhdr}
\usepackage{titlesec}
\usepackage{tcolorbox}
\usepackage{enumitem}
\usepackage{float}
\usepackage{caption}
\usepackage{microtype}
\usepackage{parskip}
\usepackage{palatino}              % Serif più elegante di Computer Modern
\usepackage{tabularx}              % Tabelle con colonne auto-adattive
\usepackage{tikz}      
\usetikzlibrary{shapes.geometric, arrows.meta, positioning}

\lstdefinestyle{bashbox}{
  language=bash,
  basicstyle=\ttfamily\footnotesize,
  columns=fullflexible,
  breaklines=true,
  breakatwhitespace=true,
  showstringspaces=false,
  frame=none,
  xleftmargin=0pt,
  xrightmargin=0pt,
  aboveskip=6pt,
  belowskip=6pt
}
\usepackage{pdfpages}              % Inclusione pagine PDF

\geometry{
  top=2.5cm, bottom=2.5cm,
  left=3cm,  right=3cm
}

% --- Liste compatte ----------------------------------------------------------
\setlist[itemize]{itemsep=3pt, parsep=2pt, topsep=4pt}
\setlist[enumerate]{itemsep=3pt, parsep=2pt, topsep=4pt}

% --- Colors ------------------------------------------------------------------
\definecolor{agentblue}{RGB}{30,80,160}
\definecolor{agentdark}{RGB}{18,50,100}
\definecolor{lightgray}{RGB}{245,245,245}
\definecolor{darkgray}{RGB}{80,80,80}
\definecolor{codebg}{RGB}{248,248,248}
\definecolor{bordercolor}{RGB}{200,200,200}
\definecolor{boxheader}{RGB}{50,50,50}
\definecolor{chaptercolor}{RGB}{25,65,140}
\definecolor{rulecolor}{RGB}{180,200,230}
\definecolor{tableheader}{RGB}{235,242,252}

% --- Hyperref ----------------------------------------------------------------
\hypersetup{
  colorlinks=true,
  linkcolor=agentblue,
  urlcolor=agentblue,
  citecolor=agentblue,
  pdftitle={Agent 2 - Consistency \& Conflict Analyzer},
  pdfauthor={Intelligent Domain Architect Team}
}

% --- tcolorbox styles --------------------------------------------------------
\tcbuselibrary{skins,breakable}

\newtcolorbox{technicalbox}[1]{
  enhanced,
  title={#1},
  colback=lightgray,
  colframe=agentblue!70!black,
  colbacktitle=agentblue,
  coltitle=white,
  fonttitle=\small\bfseries\sffamily,
  breakable,
  boxrule=0.6pt,
  arc=3pt,
  left=8pt,
  right=8pt,
  top=6pt,
  bottom=6pt,
  shadow={1pt}{-1pt}{0pt}{black!15}
}

% --- Listings ----------------------------------------------------------------
\lstset{
  basicstyle=\ttfamily\footnotesize,
  backgroundcolor=\color{codebg},
  frame=single,
  framesep=6pt,
  xleftmargin=8pt,
  xrightmargin=8pt,
  rulecolor=\color{agentblue!30},
  breaklines=true,
  breakatwhitespace=true,
  showstringspaces=false,
  tabsize=2,
  captionpos=b,
  numbers=left,
  numberstyle=\tiny\color{darkgray},
  numbersep=8pt,
  keywordstyle=\color{agentblue}\bfseries,
  commentstyle=\color{darkgray}\itshape,
  stringstyle=\color{agentblue!60!black}
}

% --- Captions ----------------------------------------------------------------
\captionsetup{
  font=small,
  labelfont={bf,color=agentblue},
  textfont={it},
  skip=8pt
}

% --- Headers/Footers ---------------------------------------------------------
\pagestyle{fancy}
\setlength{\headheight}{16pt}
\fancyhf{}
\fancyhead[L]{\small\textcolor{agentblue}{\textit{Agent 2 -- Consistency \& Conflict Analyzer}}}
\fancyhead[R]{\small\textcolor{darkgray}{\textit{Intelligent Domain Architect}}}
\fancyfoot[C]{\textcolor{darkgray}{\thepage}}
\renewcommand{\headrulewidth}{0.4pt}
\renewcommand{\headrule}{\hbox to\headwidth{\color{rulecolor}\leaders\hrule height \headrulewidth\hfill}}

% --- Chapter/Section titles --------------------------------------------------
\titleformat{\chapter}[display]
  {\normalfont\huge\bfseries\color{chaptercolor}}
  {\tikz[remember picture,overlay]\node[fill=agentblue,text=white,
    minimum width=2.5cm,minimum height=2.5cm,font=\fontsize{50}{60}\selectfont\bfseries,
    inner sep=0pt] at ([xshift=1.25cm,yshift=-1.25cm]current page.north west|-0,0)
    {\thechapter}; \hspace{2.8cm}\textcolor{darkgray}{\Large\MakeUppercase{\chaptertitlename}}}
  {10pt}{\Huge}
  [\vspace{-5pt}{\color{rulecolor}\titlerule[2pt]}]

\titleformat{\section}
  {\normalfont\Large\bfseries\color{agentblue}}{\thesection}{1em}{}
  [\vspace{-5pt}{\color{rulecolor}\titlerule[0.6pt]}]

\titleformat{\subsection}
  {\normalfont\large\bfseries\color{agentdark}}{\thesubsection}{1em}{}

% Riduce lo spazio bianco sopra i titoli di capitolo
\titlespacing*{\chapter}{0pt}{-10pt}{30pt}

% -----------------------------------------------------------------------------
\raggedbottom  % Evita troppo spazio verticale tra paragrafi

\begin{document}

% --- Title Page --------------------------------------------------------------
\begin{titlepage}
  \begin{tikzpicture}[remember picture,overlay]
    % Barra laterale sinistra
    \fill[agentblue] (current page.north west) rectangle
      ([xshift=1.2cm]current page.south west);
    % Riga decorativa in alto
    \fill[agentblue!80] ([yshift=-3cm]current page.north west) rectangle
      ([yshift=-3.08cm]current page.north east);
    % Riga decorativa in basso
    \fill[agentblue!80] ([yshift=3cm]current page.south west) rectangle
      ([yshift=3.08cm]current page.south east);
  \end{tikzpicture}

  \centering
  \vspace*{0.8cm}
  {\large\textbf{\textcolor{agentblue}{Università Campus Bio-Medico di Roma}}\\[4pt]}
  {\normalsize Facoltà di Ingegneria\\[2pt]}
  {\normalsize Corso di Laurea in Ingegneria dei Sistemi Intelligenti\\}
\vspace{3cm}
  {\color{rulecolor}\rule{0.6\textwidth}{1.5pt}}\\[0.8cm]
  {\Huge\textbf{\textcolor{agentblue}{Agent 2}}}\\[0.5cm]
  {\LARGE\textbf{\textcolor{agentdark}{Consistency \& Conflict Analyzer}}}\\[0.5cm]
  {\color{rulecolor}\rule{0.6\textwidth}{1.5pt}}\\[1cm]
  {\large Un microservizio intelligente per la validazione semantica\\
  di modelli Domain-Driven Design nella pipeline\\
  \textit{Intelligent Domain Architect}}

  \vfill

  \begin{minipage}[t]{0.45\textwidth}
    {\small\textcolor{darkgray}{\textit{Studenti:}}}\\[4pt]
    Claudio \textbf{Dragotta}\\
    Matteo \textbf{Grande}
  \end{minipage}
  \hfill
  \begin{minipage}[t]{0.45\textwidth}
    \raggedleft
    {\small\textcolor{darkgray}{\textit{Docenti:}}}\\[4pt]
    Ing.\ Floriano Caprio \\
    Dott.\ Ing. Ermanno Cordelli
  \end{minipage}

  \vspace{1.5cm}
  {\color{rulecolor}\rule{0.4\textwidth}{0.8pt}}\\[6pt]
  {\normalsize Anno Accademico 2025/2026}
\end{titlepage}

% --- ToC ---------------------------------------------------------------------
\tableofcontents
\clearpage

% --- List of Figures ---------------------------------------------------------
\listoffigures
\clearpage

% --- List of Tables ----------------------------------------------------------
\listoftables
\clearpage

% =============================================================================
\chapter{Executive Summary}
% =============================================================================

\section{Visione Strategica}

L'\textbf{Agent 2 -- Consistency \& Conflict Analyzer} è il secondo nodo della pipeline
\textit{Intelligent Domain Architect} (Agent~1 $\rightarrow$ \textbf{Agent~2} $\rightarrow$ Agent~3).
Il suo scopo è trasformare un \textit{domain model} grezzo prodotto dall'Agent~1
in un artefatto validato, coerente e pronto per la generazione automatica di specifiche
microservizi da parte dell'Agent~3.

Il sistema rileva in modo sistematico le incoerenze semantiche, i conflitti tra bounded
context e le violazioni delle regole DDD, ponendo l'utente al centro di un ciclo
iterativo di miglioramento guidato da domande contestuali.

\section{Value Proposition}

A differenza di un semplice linter sintattico, Agent~2 ragiona sul \textit{significato} del modello
attraverso tre pilastri:

\begin{itemize}
  \item \textbf{Analisi Semantica Profonda:} L'embedding \texttt{all-MiniLM-L6-v2}
        calcola la similarità vettoriale tra entità di bounded context diversi,
        rilevando sovrapposizioni (\textit{Entity Overlap}) e ambiguità terminologiche
        invisibili a un'analisi puramente sintattica.

  \item \textbf{Governance DDD Certificata:} 22 regole DDD strutturate in 7 categorie
        (classificazione dominio, bounded context, entità, ownership, requisiti, eventi,
        pattern comunicazione) vengono applicate sistematicamente, con l'LLM Qwen2.5~14B
        che approfondisce i casi complessi.

  \item \textbf{Raffinamento Iterativo Interattivo:} Il workflow n8n gestisce un loop
        adattivo in cui le risposte dell'utente vengono interpretate dall'LLM e
        applicate strutturalmente al modello prima di ogni nuova iterazione,
        convergendo verso lo stato \texttt{VALID}.
\end{itemize}
\section{Sintesi Tecnica}

L'architettura si articola in due modalità operative:

\begin{enumerate}
  \item \textbf{Modalità REST interattiva (default):} \texttt{agent2-api} espone
        12 endpoint FastAPI; n8n orchestra il loop utente tramite webhook e form HTML
        dinamici.
  \item \textbf{Modalità Kafka asincrona (predisposta):} \texttt{agent2-consumer}
        ascolta il topic \texttt{agent1-output}, processa i modelli in parallelo con
        fino a $N$ repliche e pubblica i risultati su \texttt{agent2-output} per Agent~3.
\end{enumerate}

Entrambe le modalità condividono la stessa immagine Docker e lo stesso motore di analisi.
Ollama con Qwen2.5~14B gira sull'host con accesso diretto alla GPU (NVIDIA RTX~4070,
8~GB VRAM), raggiunto dai container via \texttt{host.docker.internal:11434}.

\section{Posizionamento nella Pipeline}

Agent~2 occupa una posizione strategica come \textit{quality gate} obbligatorio tra
la fase di analisi del dominio (Agent~1) e la fase di generazione delle specifiche
tecniche (Agent~3). Questa collocazione garantisce che nessun modello DDD difettoso
raggiunga la fase di implementazione, analogamente a un gate CI/CD che blocca i
commit non conformi prima del deployment in produzione.

\section{Risultati Chiave Misurabili}

\begin{table}[H]
\centering
\renewcommand{\arraystretch}{1.3}
\begin{tabularx}{\textwidth}{l l X}
\toprule
\textbf{Metrica} & \textbf{Valore} & \textbf{Condizione} \\
\midrule
Tipologie di issue rilevate & 7 categorie & Copertura completa DDD \\
Regole DDD applicate & 22 regole & 7 categorie knowledge base \\
Tempo analisi (no LLM) & $<$5s & Modalità \texttt{use\_llm=false} \\
Timeout analisi (con LLM) & 180s max & Configurabile via .env \\
Fallback automatici & 4 livelli & LLM $\rightarrow$ euristica $\rightarrow$ template \\
Iterazioni stimate & 1--7 & Dipende da complessità modello \\
\bottomrule
\end{tabularx}
\caption{Metriche chiave di Agent~2}
\label{tab:exec-metrics}
\end{table}


% =============================================================================
\chapter{Fondamenti: Domain-Driven Design e Architetture Event-Driven}
% =============================================================================

Il presente capitolo introduce i due pilastri teorici su cui si fonda l'intera
pipeline \textit{Intelligent Domain Architect}: il \textbf{Domain-Driven Design}
(DDD) e le \textbf{architetture event-driven}. La comprensione di questi paradigmi
è essenziale per apprezzare le scelte architetturali, le regole di validazione e
i meccanismi di comunicazione tra agenti descritti nei capitoli successivi.

% -----------------------------------------------------------------------------
\section{Domain-Driven Design}
% -----------------------------------------------------------------------------

Il Domain-Driven Design è un approccio alla progettazione del software introdotto
da Eric Evans nel 2003 con il testo \textit{Domain-Driven Design: Tackling Complexity
in the Heart of Software}. L'idea centrale è che la complessità di un sistema
software risiede nel dominio di business, non nella tecnologia: di conseguenza,
il modello del dominio deve guidare ogni decisione progettuale.

\subsection{Strategic Design}

Lo \textit{Strategic Design} opera a livello macro e si occupa di suddividere un
dominio complesso in sotto-domini gestibili, ciascuno con confini ben definiti.

\begin{description}[style=nextline, leftmargin=1.5cm]
  \item[\textsf{Dominio e Sotto-domini}]
    Un \textbf{dominio} è lo spazio del problema che il software deve risolvere.
    Viene partizionato in sotto-domini classificati per rilevanza strategica:
    \begin{itemize}
      \item \textbf{Core Domain} -- il differenziatore competitivo dell'organizzazione,
            dove si concentra il massimo investimento progettuale.
      \item \textbf{Supporting Domain} -- funzionalità necessarie che supportano il
            Core, ma che non costituiscono vantaggio competitivo diretto.
      \item \textbf{Generic Domain} -- problemi già risolti dal mercato
            (autenticazione, logging, pagamenti), spesso delegabili a soluzioni
            esterne o librerie consolidate.
    \end{itemize}

  \item[\textsf{Bounded Context}]
    Un \textbf{bounded context} è il confine esplicito entro cui un determinato
    modello di dominio è valido e coerente. All'interno di un bounded context,
    ogni termine ha un significato univoco; al di fuori, lo stesso termine può
    assumere significati diversi. Ad esempio, il concetto di \texttt{Product} nel
    contesto \textit{Catalog} descrive attributi commerciali (nome, prezzo,
    descrizione), mentre nel contesto \textit{Warehouse} si riferisce a giacenza,
    posizione a scaffale e lotto di produzione.

    La corretta identificazione dei bounded context è il prerequisito per una
    architettura a microservizi coerente: ogni microservizio dovrebbe corrispondere
    a un bounded context, evitando accoppiamenti impliciti.

  \item[\textsf{Ubiquitous Language}]
    All'interno di ogni bounded context, sviluppatori ed esperti di dominio
    condividono un \textbf{linguaggio ubiquo}: un vocabolario preciso e non ambiguo
    che viene usato nel codice, nella documentazione, nelle conversazioni e nei
    test. La violazione dell'ubiquitous language -- ad esempio usare sinonimi
    informali o termini tecnici non condivisi -- è una delle cause più frequenti
    di disallineamento tra requisiti e implementazione.

  \item[\textsf{Context Map}]
    La \textbf{Context Map} è la rappresentazione delle relazioni tra bounded
    context. Evans definisce diversi pattern di integrazione, tra cui:
    \begin{itemize}
      \item \textbf{Shared Kernel} -- due contesti condividono una porzione di
            modello, con responsabilità congiunta sulle modifiche.
      \item \textbf{Customer-Supplier} -- un contesto (supplier) fornisce dati
            a un altro (customer), con contratti espliciti sulle interfacce.
      \item \textbf{Conformist} -- il contesto a valle adotta il modello del
            contesto a monte senza negoziare, subendo le modifiche imposte.
      \item \textbf{Anticorruption Layer (ACL)} -- il contesto a valle costruisce
            uno strato di traduzione che protegge il proprio modello dalle
            contaminazioni del modello a monte.
      \item \textbf{Open Host Service} -- il contesto a monte espone un servizio
            ben definito (es.\ API REST) che altri contesti possono consumare.
      \item \textbf{Published Language} -- un formato di scambio condiviso
            (es.\ JSON Schema, Protobuf, Avro) definisce il contratto di
            comunicazione tra contesti.
      \item \textbf{Separate Ways} -- nessuna integrazione; i contesti operano
            in completa autonomia, utile per domini di supporto non correlati.
    \end{itemize}
\end{description}

\subsection{Tactical Design}

Il \textit{Tactical Design} opera a livello micro, all'interno di un singolo
bounded context, e fornisce i building block per modellare il dominio nel codice.

\begin{description}[style=nextline, leftmargin=1.5cm]
  \item[\textsf{Entity}]
    Un oggetto definito dalla sua \textbf{identità} piuttosto che dai suoi
    attributi. Due entità con gli stessi attributi ma identità diverse sono oggetti
    distinti (es.\ due clienti con lo stesso nome ma \texttt{id} diverso).

  \item[\textsf{Value Object}]
    Un oggetto \textbf{immutabile} definito esclusivamente dai suoi attributi, privo
    di identità propria. Due value object con gli stessi attributi sono considerati
    equivalenti (es.\ un indirizzo, un importo monetario, un intervallo di date).

  \item[\textsf{Aggregate}]
    Un cluster di entity e value object trattato come un'\textbf{unità transazionale}
    unica. L'accesso dall'esterno avviene esclusivamente attraverso l'\textbf{Aggregate
    Root}, che garantisce le invarianti di business. Ad esempio, un \texttt{Order}
    (aggregate root) contiene \texttt{OrderLine} (entity) e \texttt{Money} (value
    object); nessun componente esterno modifica direttamente una \texttt{OrderLine}
    senza passare da \texttt{Order}.

  \item[\textsf{Domain Event}]
    Un fatto significativo che è accaduto nel dominio, espresso al \textbf{passato}
    (es.\ \texttt{OrderPlaced}, \texttt{PaymentConfirmed}). I domain event sono il
    meccanismo naturale attraverso cui i bounded context comunicano in modo
    disaccoppiato: un contesto emette un evento, altri contesti interessati lo
    consumano e reagiscono secondo la propria logica.

  \item[\textsf{Repository}]
    L'astrazione che fornisce accesso agli aggregati come se fossero in una
    collezione in memoria, nascondendo i dettagli di persistenza.

  \item[\textsf{Domain Service}]
    Logica di dominio che non appartiene naturalmente a nessuna entity o value
    object, ma che opera su più aggregati o richiede collaborazione tra essi
    (es.\ un servizio di calcolo tariffe che coinvolge \texttt{Customer},
    \texttt{Product} e \texttt{PricingPolicy}).
\end{description}

\begin{technicalbox}{Perché il DDD è rilevante per Agent~2}
Agent~2 valida un domain model verificando il rispetto di questi principi:
controlla che ogni bounded context abbia un linguaggio coerente, che le entità
non siano duplicate tra contesti, che gli aggregati rispettino le invarianti,
che gli eventi seguano la naming convention al passato e che i pattern di
comunicazione nella context map siano compatibili. Le 22~regole della knowledge
base di Agent~2 sono una formalizzazione operativa dei principi DDD descritti
in questa sezione.
\end{technicalbox}

% -----------------------------------------------------------------------------
\section{Architetture Event-Driven}
% -----------------------------------------------------------------------------

Un'\textbf{architettura event-driven} (EDA) è un paradigma in cui i componenti del
sistema comunicano attraverso la produzione, il rilevamento e il consumo di
\textit{eventi}. A differenza delle architetture request-response tradizionali, in
cui un componente invoca esplicitamente un altro e attende la risposta, nell'EDA
il produttore emette un evento senza conoscere i consumatori: questo
\textit{disaccoppiamento} è il vantaggio fondamentale del paradigma.

\subsection{Concetti Chiave}

\begin{description}[style=nextline, leftmargin=1.5cm]
  \item[\textsf{Evento}]
    Un record immutabile che descrive un fatto accaduto nel sistema.
    Contiene tipicamente un tipo (es.\ \texttt{OrderPlaced}), un timestamp, un
    payload con i dati rilevanti e metadati di correlazione.

  \item[\textsf{Producer e Consumer}]
    Il \textbf{producer} è il componente che emette l'evento; il \textbf{consumer}
    è il componente che lo riceve e reagisce. Un singolo evento può avere zero,
    uno o molti consumer, abilitando pattern \textit{fan-out} senza modifiche al
    producer.

  \item[\textsf{Broker / Message Bus}]
    L'infrastruttura che media tra producer e consumer. Il broker riceve gli eventi,
    li persiste (garantendo durabilità) e li consegna ai consumer registrati.
    Esempi diffusi includono Apache Kafka, RabbitMQ e Amazon SNS/SQS.

  \item[\textsf{Topic e Partizione}]
    In broker come Apache Kafka, gli eventi vengono pubblicati su \textbf{topic}
    logici. Ogni topic è suddiviso in \textbf{partizioni} che consentono il
    parallelismo: più consumer possono leggere partizioni diverse
    contemporaneamente, scalando orizzontalmente il throughput.
\end{description}

\subsection{Pattern di Comunicazione}

Le architetture event-driven si articolano in diversi pattern, ciascuno con
caratteristiche e trade-off specifici:

\begin{description}[style=nextline, leftmargin=1.5cm]
  \item[\textsf{Event Notification}]
    Il producer emette un evento minimale (es.\ \texttt{OrderPlaced\{orderId:
    "123"\}}) che segnala che qualcosa è accaduto, senza includere tutti i dati.
    I consumer interessati effettuano una callback al producer per recuperare i
    dettagli. Questo pattern minimizza la dimensione dei messaggi ma introduce
    un accoppiamento temporale.

  \item[\textsf{Event-Carried State Transfer}]
    L'evento contiene tutti i dati necessari affinché i consumer possano operare
    senza callback (es.\ \texttt{OrderPlaced\{orderId, items, total, customer\}}).
    Elimina l'accoppiamento temporale a costo di messaggi più pesanti e della
    necessità di gestire la consistenza eventuale dei dati replicati.

  \item[\textsf{Event Sourcing}]
    Lo stato di un'entità non viene memorizzato come snapshot, ma come sequenza
    ordinata di tutti gli eventi che lo hanno prodotto. Lo stato corrente si
    ottiene riproducendo (\textit{replaying}) la sequenza. Questo pattern offre
    un audit trail completo e la possibilità di ricostruire lo stato a qualsiasi
    punto nel tempo, ma aumenta la complessità di query e gestione.

  \item[\textsf{CQRS (Command Query Responsibility Segregation)}]
    Separa il modello di \textit{scrittura} (command) dal modello di
    \textit{lettura} (query). I comandi producono eventi che aggiornano il modello
    di scrittura; proiezioni asincrone materializzano viste ottimizzate per la
    lettura. CQRS si combina frequentemente con Event Sourcing per ottenere
    massima scalabilità e flessibilità.
\end{description}

\subsection{Gestione dei Flussi di Business: Orchestrazione vs Coreografia}

Quando un processo di business coinvolge più servizi (es.\ creazione ordine
$\rightarrow$ pagamento $\rightarrow$ aggiornamento inventario $\rightarrow$
notifica), esistono due strategie per coordinare il flusso:

\begin{description}[style=nextline, leftmargin=1.5cm]
  \item[\textsf{Orchestrazione}]
    Un \textbf{orchestratore centrale} conosce l'intero flusso di lavoro: riceve
    un evento iniziale, invoca i servizi nella sequenza corretta e gestisce gli
    errori. La logica di business è centralizzata, il che semplifica la comprensione
    del flusso ma introduce un \textit{single point of failure} e un forte
    accoppiamento (l'orchestratore deve conoscere tutti i servizi).

  \item[\textsf{Coreografia}]
    Nessun controllore centrale: ogni servizio è \textbf{indipendente} e reagisce
    agli eventi pubblicati dagli altri. Ad esempio, \texttt{OrderService} pubblica
    \texttt{OrderCreated}; \texttt{PaymentService} lo ascolta, processa il
    pagamento e pubblica \texttt{PaymentProcessed}; \texttt{InventoryService}
    ascolta \texttt{PaymentProcessed} e aggiorna la giacenza. Il sistema è molto
    disaccoppiato e resiliente, ma il flusso complessivo è distribuito e più
    difficile da monitorare.
\end{description}

\begin{technicalbox}{Orchestrazione vs Coreografia: Quando Usare Quale}
\textbf{Orchestrazione:} flussi complessi con logica di compensazione esplicita
(es.\ saga pattern), dove il monitoraggio centralizzato è prioritario.
Se un servizio muore, l'orchestratore se ne accorge immediatamente e può
eseguire retry o segnare l'operazione come fallita.

\textbf{Coreografia:} sistemi altamente scalabili dove il disaccoppiamento è
prioritario. Se un servizio muore, i messaggi non consumati restano nel broker
e verranno elaborati al ripristino. Ogni consumer è potenzialmente anche un
producer, creando una catena reattiva di eventi.
\end{technicalbox}

\subsection{Vantaggi e Sfide}

\begin{table}[H]
\centering
\renewcommand{\arraystretch}{1.3}
\begin{tabularx}{\textwidth}{>{\bfseries}l X}
\toprule
\textbf{Vantaggi} & \textbf{Descrizione} \\
\midrule
Disaccoppiamento   & Producer e consumer evolvono indipendentemente; nuovi
                     consumer possono essere aggiunti senza modificare i producer. \\
Scalabilità        & Il parallelismo a livello di partizioni consente di scalare
                     orizzontalmente il throughput di elaborazione. \\
Resilienza         & Il broker persiste gli eventi; se un consumer è
                     temporaneamente non disponibile, elaborerà i messaggi al
                     ripristino. \\
Estensibilità      & Aggiungere nuove reazioni a un evento non richiede
                     modifiche ai componenti esistenti (principio Open-Closed). \\
\bottomrule
\end{tabularx}
\caption{Vantaggi delle architetture event-driven}
\label{tab:eda-vantaggi}
\end{table}

\begin{table}[H]
\centering
\renewcommand{\arraystretch}{1.3}
\begin{tabularx}{\textwidth}{>{\bfseries}l X}
\toprule
\textbf{Sfide} & \textbf{Descrizione} \\
\midrule
Consistenza eventuale & I dati tra servizi non sono immediatamente allineati;
                        occorre progettare per la \textit{eventual consistency}. \\
Osservabilità         & Tracciare il flusso di un evento attraverso molteplici
                        consumer richiede strumenti di tracing distribuito
                        (es.\ correlation ID). \\
Ordinamento           & Garantire l'ordine di elaborazione degli eventi è
                        possibile solo all'interno di una singola partizione. \\
Complessità operativa & La gestione di un broker (configurazione, monitoraggio,
                        retention) aggiunge overhead infrastrutturale. \\
\bottomrule
\end{tabularx}
\caption{Sfide delle architetture event-driven}
\label{tab:eda-sfide}
\end{table}

\subsection{Apache Kafka come Infrastruttura Event-Driven}

Apache Kafka è la piattaforma di event streaming più diffusa per architetture
EDA ad alta disponibilità. A differenza di una coda di messaggi tradizionale,
Kafka è fondamentalmente un \textbf{log distribuito}: un file append-only,
immutabile e ordinato. Questa scelta progettuale elimina le operazioni più
costose di un database (letture e scritture casuali), rendendo Kafka
estremamente veloce e affidabile.

\begin{description}[style=nextline, leftmargin=1.5cm]
  \item[\textsf{Log Immutabile}]
    Ogni evento scritto su Kafka viene aggiunto in coda al log e non può essere
    modificato né cancellato (almeno fino alla scadenza della retention policy).
    La scrittura è un'operazione $O(1)$; la lettura avviene per offset sequenziale,
    altrettanto veloce.

  \item[\textsf{Partizioni e Scalabilità}]
    Un topic non è un singolo file ma viene suddiviso in $N$
    \textbf{partizioni}, distribuite su più server (\textit{broker}) nel cluster.
    Ogni partizione è un log ordinato indipendente. Più partizioni consentono a
    più producer e consumer di operare in parallelo, scalando orizzontalmente
    il throughput.

  \item[\textsf{Replicazione Leader/Follower}]
    Per ogni partizione, Kafka mantiene un \textbf{Leader} (che riceve tutte le
    scritture) e uno o più \textbf{Follower} (che replicano i dati dal Leader).
    Se il Leader diventa non disponibile, un Follower aggiornato viene eletto
    automaticamente come nuovo Leader, garantendo \textbf{tolleranza ai guasti}
    senza perdita di dati.

  \item[\textsf{Consumer Group}]
    Più istanze dello stesso consumer possono essere raggruppate in un
    \textbf{consumer group}. Kafka assegna automaticamente le partizioni ai
    consumer del gruppo: se il topic ha 4~partizioni e il gruppo ha 2~consumer,
    ciascuno leggerà 2~partizioni. Se un consumer muore, le sue partizioni
    vengono riassegnate agli altri (\textit{rebalancing}), garantendo
    parallelismo e resilienza anche in lettura.

  \item[\textsf{Acknowledgments (ACKs)}]
    Il livello di garanzia sulla consegna è configurabile tramite il parametro
    \texttt{acks} del producer:
    \begin{itemize}
      \item \texttt{acks=0} -- il producer non attende conferma (massima velocità,
            rischio di perdita).
      \item \texttt{acks=1} -- il producer attende la conferma del solo Leader
            (buon compromesso velocità/sicurezza).
      \item \texttt{acks=all} -- il producer attende la conferma di Leader e tutti
            i Follower (massima sicurezza, latenza più alta).
    \end{itemize}
\end{description}

\subsection{EDA Stateful vs Stateless e il Wrapper Sincrono}

Le architetture event-driven si distinguono anche per la gestione dello stato:

\begin{description}[style=nextline, leftmargin=1.5cm]
  \item[\textsf{Stateless EDA}]
    Gli eventi sono autonomi e non correlati tra loro. Ogni evento contiene
    tutte le informazioni necessarie per essere elaborato indipendentemente
    (es.\ un evento di telemetria GPS).

  \item[\textsf{Stateful EDA}]
    Gli eventi sono correlati e il loro significato dipende dalla sequenza o
    dallo stato accumulato (es.\ una saga di ordine che attraversa gli stati
    \texttt{Created} $\rightarrow$ \texttt{Paid} $\rightarrow$ \texttt{Shipped}).
    La gestione dello stato richiede meccanismi come Event Sourcing o proiezioni
    materializzate.
\end{description}

Un problema ricorrente è l'interfacciamento con \textbf{client sincroni}
(es.\ un'applicazione mobile o un browser) che necessitano di una risposta
immediata, mentre l'EDA è intrinsecamente asincrona. La soluzione è il
\textbf{Wrapper sincrono}: un componente che (1) riceve la richiesta REST dal
client, (2) genera un evento asincrono e lo pubblica sul broker, (3) si
iscrive a un topic di risposta e attende, (4) restituisce la risposta HTTP
al client quando l'evento di completamento arriva. Questo pattern
``disintermedia'' la comunicazione sincrona/asincrona senza sacrificare il
disaccoppiamento interno.

% -----------------------------------------------------------------------------
\section{DDD, Microservizi e Cloud: la Triade Architetturale}
% -----------------------------------------------------------------------------

DDD, architetture event-driven e cloud computing formano una triade
sinergica per la costruzione di sistemi moderni:

\begin{itemize}
  \item \textbf{Microservizi come scomposizione fisica:} ogni bounded context
        del DDD diventa un microservizio autonomo, con il proprio database,
        la propria logica di dominio e il proprio ciclo di deployment.

  \item \textbf{EDA come comunicazione disaccoppiata:} i domain event
        implementano la comunicazione tra bounded context attraverso un
        broker (Kafka), eliminando l'accoppiamento sincrono tipico delle
        architetture REST point-to-point.

  \item \textbf{Cloud come infrastruttura elastica:} ogni microservizio
        (bounded context) può essere scalato indipendentemente in base al
        carico. La context map del DDD diventa la mappa delle interazioni
        tra servizi cloud, e l'Anticorruption Layer protegge il modello di
        dominio dalle API di servizi esterni.
\end{itemize}

Questa triade è alla base della pipeline \textit{Intelligent Domain Architect}:
Agent~1, Agent~2 e Agent~3 sono microservizi indipendenti che comunicano via
Kafka (EDA), ciascuno con un bounded context ben definito (DDD), containerizzati
e scalabili tramite Docker Compose (infrastruttura).

% -----------------------------------------------------------------------------
\section{DDD e Event-Driven nei Sistemi Multi-Agente}
% -----------------------------------------------------------------------------

La pipeline \textit{Intelligent Domain Architect} rappresenta un caso concreto in
cui DDD e architetture event-driven si integrano in un sistema multi-agente:

\begin{enumerate}
  \item \textbf{Il dominio come artefatto centrale.} L'intero flusso ruota attorno
        a un \textit{domain model} strutturato secondo i principi DDD: bounded
        context, entità, aggregati, eventi e pattern di comunicazione. Il modello
        non è un diagramma statico, ma un artefatto vivo che viene analizzato,
        validato e raffinato iterativamente dagli agenti.

  \item \textbf{Comunicazione event-driven tra agenti.} Gli agenti della pipeline
        comunicano tramite Apache Kafka: Agent~1 pubblica il domain model sul topic
        \texttt{agent1-output}, Agent~2 lo consuma, esegue l'analisi e pubblica il
        risultato su \texttt{agent2-output} per Agent~3. Questo disaccoppiamento
        consente a ciascun agente di evolvere, scalare e fallire indipendentemente
        dagli altri.

  \item \textbf{Validazione DDD come quality gate.} Agent~2 formalizza i principi
        DDD in controlli automatizzati (cfr.\ Capitolo~6). Ogni violazione produce
        un \textit{issue} tipizzato che viene restituito all'utente con domande
        contestuali per guidare la correzione.

  \item \textbf{Consistenza eventuale e iterazione.} Come in ogni sistema
        event-driven, la consistenza del modello non è garantita al primo
        passaggio. Agent~2 adotta un approccio iterativo: ad ogni ciclo, le
        risposte dell'utente vengono interpretate dall'LLM e applicate
        strutturalmente al modello, riducendo progressivamente le issue fino al
        raggiungimento dello stato \texttt{VALID}.
\end{enumerate}

\begin{technicalbox}{Dalla Teoria alla Pratica}
I capitoli successivi mostreranno come ciascun concetto introdotto in questa sezione
si traduce in componenti concreti del sistema:
\begin{itemize}
  \item I \textbf{bounded context} e le \textbf{entità} DDD diventano le strutture
        dati JSON analizzate dalla pipeline (Capitolo~8).
  \item Le \textbf{regole DDD} diventano la knowledge base del
        \texttt{ConflictDetector} (Capitolo~6).
  \item I \textbf{domain event} diventano i messaggi Kafka tra agenti (Capitolo~7).
  \item I \textbf{pattern della context map} diventano i controlli di compatibilità
        nella comunicazione tra bounded context (Capitolo~7).
\end{itemize}
\end{technicalbox}

% =============================================================================
\chapter{Analisi dello Scenario e Business Case}
% =============================================================================

\section{Razionale del Progetto}

La progettazione di sistemi software complessi richiede la definizione di un \textit{domain
model} coerente prima di avviare lo sviluppo dei microservizi. Nella pratica, i modelli
prodotti nella fase di analisi -- anche quando generati automaticamente da Agent~1 --
contengono regolarmente incoerenze strutturali: entità duplicate tra bounded context,
requisiti contraddittori, eventi con naming scorretto, pattern di comunicazione
incompatibili.

Senza un passaggio di validazione formale, questi problemi si propagano fino all'Agent~3,
che genererebbe specifiche difettose, costose da correggere nelle fasi successive.

\section{Analisi dei Pain Points}

Lo scenario attuale, in assenza di un validatore semantico, presenta criticità strutturali:

\begin{itemize}
  \item \textbf{Entity Overlap silenzioso:} La stessa entità (es.\ \texttt{Product})
        compare in più bounded context con significati leggermente diversi, violando
        il principio di \textit{Single Source of Truth} senza che nessuno strumento
        convenzionale lo rilevi.

  \item \textbf{Conflitti di Requisiti:} Requisiti funzionali di domini diversi si
        contraddicono (es.\ ``immutabile dopo la creazione'' vs.\ ``modificabile fino
        alla spedizione'') generando ambiguità che emergono solo in fase di sviluppo.

  \item \textbf{Violazioni dell'Ubiquitous Language:} Termini come
        \texttt{CreateOrder} (anziché \texttt{OrderCreated}) o domini di supporto
        classificati erroneamente come Core distorcono le responsabilità architetturali.

  \item \textbf{Pattern Incompatibili:} Un dominio richiede consistenza forte con
        comunicazione sincrona, ma l'architettura eventi prevede solo canali asincroni,
        rendendo il contratto di integrazione irrealizzabile.
\end{itemize}

\section{Scenario d'Uso Concreto: E-Commerce}

Per illustrare il valore di Agent~2, consideriamo un dominio e-commerce tipico
in cui Agent~1 ha prodotto un domain model con le seguenti caratteristiche:

\begin{itemize}
\item 3 domini core: \texttt{OrderManagement}, \texttt{CatalogManagement},
        \texttt{PaymentManagement}
  \item 2 domini supporting: \texttt{NotificationService}, \texttt{UserManagement}
  \item 15 entità distribuite tra i bounded context
  \item 8 eventi di dominio
  \item 5 requisiti funzionali per dominio
\end{itemize}

In questo scenario, Agent~2 tipicamente rileva:

\begin{table}[H]
\centering
\small
\renewcommand{\arraystretch}{1.5}
\setlength{\tabcolsep}{6pt}

\begin{tabularx}{\textwidth}{
  >{\raggedright\arraybackslash}p{8.0cm}
  >{\raggedright\arraybackslash}X
}
\toprule
\textbf{Issue} & \textbf{Descrizione} \\
\midrule

\texttt{ENTITY\_OVERLAP} \textbf{(HIGH)}
& \texttt{Product} presente sia in \texttt{OrderContext} sia in \texttt{CatalogContext} \\

\texttt{REQUIREMENT\_CONFLICT} \textbf{(HIGH)}
& immutabilità dell'ordine vs.\ modificabilità del prodotto \\

\texttt{NAMING\_VIOLATION} \textbf{(MEDIUM)}
& \texttt{CreateOrder} $\rightarrow$ \texttt{OrderCreated} \\

\texttt{MISCLASSIFIED\_DOMAIN} \textbf{(MEDIUM)}
& \texttt{UserManagement} classificato come \texttt{Core} invece di \texttt{Generic} \\

\bottomrule
\end{tabularx}
\end{table}


Il loop n8n guida il business stakeholder attraverso 2 iterazioni: la prima per
risolvere l'entity overlap e il conflitto di requisiti, la seconda per confermare
le scelte e raggiungere lo stato \texttt{VALID}.

\section{Scenario d'Uso: Healthcare}

Il dominio sanitario introduce complessità semantiche specifiche che Agent~2
gestisce tramite la knowledge base estesa:

\begin{itemize}
  \item \texttt{Patient} appare in \texttt{ClinicalContext} (persona in cura),
        \texttt{RecordsContext} (titolare della cartella clinica) e
        \texttt{BillingContext} (soggetto pagante).
        L'embedding rileva la sovrapposizione con similarità $>0.85$.
  \item Il termine \texttt{Visita} ha significati diversi: in
        \texttt{ClinicalContext} indica l'appuntamento del paziente con il
        medico, mentre in \texttt{RecordsContext} indica il referto prodotto
        a seguito del consulto, generando \texttt{SEMANTIC\_AMBIGUITY}.
  \item Il file di test \texttt{case\_healthcare.json} in \texttt{input\_agent/}
        contiene questi 4 issue precostituiti per la verifica del sistema.
\end{itemize}
\section{Rischi e Impatto Organizzativo}

Il mancato rilevamento di questi problemi nella fase di modellazione eleva
significativamente il costo del ciclo di sviluppo. Agent~2 mitiga tali rischi
attraverso un processo di validazione iterativa che coinvolge attivamente il business
stakeholder, garantendo che ogni modifica al domain model sia tracciabile, motivata
e approvata prima di procedere alla specifica dei microservizi.

\begin{table}[H]
\centering
\renewcommand{\arraystretch}{1.3}
\begin{tabularx}{\textwidth}{>{\raggedright\arraybackslash}l >{\raggedright\arraybackslash}X >{\raggedright\arraybackslash}X}
\toprule
\textbf{Rischio} & \textbf{Senza Agent~2} & \textbf{Con Agent~2} \\
\midrule
Entity Overlap & Silenzioso fino all'implementazione & Rilevato e risolto prima di Agent~3 \\
Conflitti di requisiti & Scoperto durante lo sviluppo & Segnalato con domande contestuali \\
Naming violations & Propagate alle API generate & Auto-corrette dallo Step~4 \\
Decisioni non tracciate & Persa la storia delle scelte & Audit trail nei metadati del modello \\
\bottomrule
\end{tabularx}
\caption{Confronto rischi con e senza Agent~2}
\label{tab:risk-comparison}
\end{table}

% =============================================================================
\chapter{Obiettivi e Requisiti del Sistema}
% =============================================================================

\section{Obiettivi Strategici}

Agent~2 opera come microservizio autonomo di \textbf{Domain Model Governance}.
Gli obiettivi si articolano in:

\begin{itemize}
  \item \textbf{Validazione Semantica Automatizzata:} Rilevamento sistematico di tutte le
        categorie di incoerenza DDD prima che il modello raggiunga Agent~3.

  \item \textbf{Raffinamento Guidato dal Contesto:} Generazione di domande specifiche per
        nome di entità e bounded context, con risposte suggerite azionabili,
        che guidano il stakeholder verso decisioni concrete.

  \item \textbf{Convergenza Garantita:} Il loop iterativo con limite adattivo assicura
        che il sistema non si blocchi prematuramente, proseguendo fino al raggiungimento
        dello stato \texttt{VALID}.

  \item \textbf{Interoperabilità A2A:} Conformità al protocollo A2A v0.3 per la
        comunicazione strutturata con Agent~1 e Agent~3 nella pipeline.
\end{itemize}

\section{Requisiti Funzionali}

\begin{enumerate}
  \item \textbf{Pipeline di Analisi in 5 Step:} Ogni chiamata a \texttt{POST /analyze}
        esegue sequenzialmente: pre-applicazione risposte, analisi semantica,
        conflict detection, generazione domande, raffinamento modello.

  \item \textbf{Rilevamento Multi-Categoria:} Il sistema deve identificare almeno
        7 tipologie di issue: \texttt{ENTITY\_OVERLAP}, \texttt{SEMANTIC\_AMBIGUITY},
        \texttt{REQUIREMENT\_CONFLICT}, \texttt{DUPLICATE\_EVENT},
        \texttt{NAMING\_VIOLATION}, \texttt{INCOMPATIBLE\_PATTERN},
        \texttt{MISCLASSIFIED\_DOMAIN}.

  \item \textbf{LLM-first con Fallback:} La generazione delle domande utilizza
        Qwen2.5~14B come approccio primario; in caso di timeout o JSON malformato,
        attiva automaticamente i template dalla \texttt{validation\_checklist.json}.

  \item \textbf{Filtraggio Intelligente:} Gli issue relativi ad elementi già risolti
        nelle iterazioni precedenti (\texttt{\_userGuidedFixes},
        \texttt{\_ownershipDecisions}) vengono filtrati per non riproporre
        problemi già trattati.

  \item \textbf{Stima Adattiva delle Iterazioni:} Il campo
        \texttt{suggested\_max\_iterations} viene ricalcolato a ogni iterazione
        tramite LLM (con fallback euristico) per adattare dinamicamente il limite
        del loop n8n.

  \item \textbf{Doppia Modalità Operativa:} Supporto alla modalità REST interattiva
        (via n8n) e alla modalità Kafka consumer asincrona per la pipeline automatica
        Agent~1 $\rightarrow$ Agent~2 $\rightarrow$ Agent~3.
\end{enumerate}

\section{Requisiti Non Funzionali}

\begin{itemize}
  \item \textbf{Performance:} La pipeline completa con LLM deve completarsi entro
        180~secondi per timeout; la modalità \texttt{use\_llm=false} garantisce
        risposte in pochi secondi per scenari batch.

  \item \textbf{Scalabilità:} Il servizio Kafka consumer supporta scaling orizzontale
        tramite la variabile \texttt{CONSUMER\_REPLICAS} in docker-compose, con
        processamento parallelo fino a \texttt{MAX\_WORKERS=4} thread per replica.

  \item \textbf{Portabilità:} Containerizzazione completa tramite Docker Compose
        multi-file per isolamento delle dipendenze e replicabilità dell'ambiente.

  \item \textbf{Resilienza AI:} Fallback automatico da LLM a modalità euristica/template
        in caso di indisponibilità di Ollama, garantendo la continuità del servizio.

  \item \textbf{Interoperabilità:} Conformità allo standard A2A Protocol v0.3
        con Agent Card pubblicata su \texttt{/.well-known/agent.json}.
\end{itemize}

\section{Casi d'Uso Principali}

\subsection{UC-001: Analisi Iniziale di un Domain Model}

\textbf{Attore:} Business Stakeholder tramite n8n \\
\textbf{Pre-condizione:} Docker stack avviato, Ollama con Qwen2.5~14B disponibile \\
\textbf{Flusso:}
\begin{enumerate}
  \item L'utente apre \texttt{http://127.0.0.1:5678/webhook/agent2-start} nel browser.
  \item Seleziona la modalità di input (file predefinito o upload) e il domain model.
  \item n8n invia il modello a \texttt{POST /analyze} con \texttt{use\_llm=true}.
  \item Agent~2 esegue i 4 step di analisi e restituisce issue e domande.
  \item n8n renderizza le domande come form HTML con radio buttons.
  \item L'utente risponde e il ciclo riparte dalla chiamata successiva a \texttt{/analyze}.
  \item Quando \texttt{status=VALID}, n8n mostra la pagina di successo con il modello raffinato.
\end{enumerate}

\subsection{UC-002: Analisi Batch via Kafka (Predisposta)}

\textbf{Attore:} Agent~1 (automatico, nessuna interazione umana) \\
\textbf{Pre-condizione:} Stack Kafka avviato con \texttt{docker-compose.kafka.yml} \\
\textbf{Flusso:}
\begin{enumerate}
  \item Agent~1 pubblica il domain model su topic \texttt{agent1-output}.
  \item \texttt{agent2-consumer} consuma il messaggio dal topic.
  \item Esegue la pipeline di analisi completa (modalità non interattiva, \texttt{use\_llm=true}).
  \item Pubblica il modello validato e raffinato su topic \texttt{agent2-output}.
  \item Agent~3 consuma il risultato e avvia la generazione delle specifiche.
\end{enumerate}

\subsection{UC-003: Debug e Test Diretto API}

\textbf{Attore:} Team DevOps o sviluppatore \\
\textbf{Pre-condizione:} Agent~2 API avviata \\
\textbf{Flusso:}
\begin{enumerate}
  \item Accede a \texttt{http://localhost:8002/docs} (Swagger UI).
  \item Seleziona un file di test predefinito tramite \texttt{GET /source/files}.
  \item Invia l'analisi con \texttt{POST /source/analyze-by-file?filename=example\_bad.json}.
  \item Verifica il risultato JSON con gli issue attesi per lo stress test completo.
  \item Consulta i log tramite \texttt{docker logs agent2-api} per la diagnostica.
\end{enumerate}

\section{Requisiti di Sistema}

\begin{table}[H]
\centering
\renewcommand{\arraystretch}{1.3}
\begin{tabular}{llp{5.5cm}}
\toprule
\textbf{Componente} & \textbf{Minimo} & \textbf{Raccomandato} \\
\midrule
RAM sistema & 4 GB liberi & 8 GB (con Kafka: 12 GB) \\
CPU & 2 core & 4+ core \\
GPU NVIDIA & Non richiesta (CPU mode) & RTX 4070 / 8 GB VRAM \\
Storage & 10 GB & 20 GB (modelli LLM inclusi) \\
OS & Linux / macOS / Windows 10+ & Windows 11 con WSL2 o Linux \\
Docker & 24.0+ & Ultima versione stabile \\
Ollama & 0.3+ & 0.5+ con modello qwen2.5:14b \\
\bottomrule
\end{tabular}
\caption{Requisiti hardware e software di sistema}
\label{tab:system-requirements}
\end{table}

% =============================================================================
\chapter{Stakeholder e Ruoli}
% =============================================================================

Il presente capitolo identifica gli attori che interagiscono con Agent~2,
sia come utenti umani che come sistemi della pipeline.

\section{Business Stakeholder (Utente Interattivo)}

\textbf{Ruolo:} Architetto software, domain expert o product owner che partecipa
attivamente al ciclo di raffinamento del modello tramite l'interfaccia n8n.

\textbf{Responsabilità:} Rispondere alle domande contestuali generate dall'agente,
scegliere tra le risposte suggerite o fornire input libero, validare il modello
raffinato prima dell'inoltro ad Agent~3.

\begin{technicalbox}{Dettaglio Tecnico: Interazione via n8n}
Il workflow \texttt{workflow\_complete\_loop.json} presenta all'utente un form HTML
dinamico con radio buttons per le risposte suggerite e campo testo libero.
Le risposte vengono inviate via \texttt{POST /webhook/agent2-submit} e passate
all'iterazione successiva come \texttt{previous\_answers}.
\end{technicalbox}

\section{Agent 1 -- Domain Interviewer (Upstream)}

\textbf{Ruolo:} Produttore del domain model JSON che alimenta Agent~2.

\textbf{Responsabilità:} Costruire un domain model strutturato tramite intervista
al business stakeholder e pubblicarlo sul topic Kafka \texttt{agent1-output}
oppure inviarlo direttamente via \texttt{POST /a2a/message/send}.

\begin{technicalbox}{Dettaglio Tecnico: Integrazione A2A e Kafka}
In modalità REST, Agent~1 invia il modello tramite il protocollo A2A v0.3:
\texttt{POST /a2a/message/send} con un \texttt{Part} contenente il JSON del
domain model. In modalità Kafka (futura), pubblica il messaggio sul topic
\texttt{agent1-output}; l'\texttt{agent2-consumer} lo consuma automaticamente.
\end{technicalbox}

\section{Agent 3 -- Specification Generator (Downstream)}

\textbf{Ruolo:} Consumatore del domain model validato e raffinato prodotto da Agent~2.

\textbf{Responsabilità:} Ricevere il \texttt{refined\_model} in stato \texttt{VALID}
e generare le specifiche tecniche dei microservizi.

\begin{technicalbox}{Dettaglio Tecnico: Output verso Agent 3}
Il modello raffinato viene salvato su \texttt{output\_agent/\{filename\}/\{session\_id\}/iteration\_\{N\}.json}
e, nella modalità Kafka futura, pubblicato sul topic \texttt{agent2-output}.
Agent~3 consuma questo topic per avviare la fase di specifica.
\end{technicalbox}

\section{Operatore n8n (Workflow Engineer)}

\textbf{Ruolo:} Responsabile della configurazione e manutenzione del workflow n8n.

\textbf{Responsabilità:} Importare e attivare i workflow dalla directory \texttt{n8n/},
configurare gli URL degli endpoint (\texttt{http://agent2-api:8002}), monitorare
il loop interattivo e aggiornare i workflow in caso di evoluzioni dell'API.

\section{Team Tecnico e DevOps}

\textbf{Ruolo:} Sviluppatori e sistemisti incaricati della manutenzione evolutiva
del microservizio.

\textbf{Responsabilità:} Gestione del ciclo di vita Docker, aggiornamento dei
modelli LLM in Ollama, evoluzione della knowledge base DDD, integrazione Kafka
per la pipeline automatica.

\begin{technicalbox}{Dettaglio Tecnico: Infrastruttura Docker}
Lo stack Docker è descritto in dettaglio nel Capitolo~10 (Deployment e Operatività).
In sintesi: \texttt{docker-compose.yml} avvia i 2 container core,
\texttt{docker-compose.kafka.yml} aggiunge i 4 container per la pipeline Kafka.
\end{technicalbox}

\section{Matrice di Tracciabilità Endpoint/Attore}

\begin{table}[H]
\centering
\renewcommand{\arraystretch}{1.3}
\begin{tabular}{lcccc}
\toprule
\textbf{Attore} & \textbf{/analyze} & \textbf{/a2a/*} & \textbf{/source/*} & \textbf{/health} \\
\midrule
Business Stakeholder (n8n) & Sì & No  & Sì  & No  \\
Agent 1 (A2A)              & No & Sì  & No  & No  \\
Agent 3 (Kafka)            & No & No  & No  & No  \\
DevOps / Debug             & Sì & Sì  & Sì  & Sì  \\
\bottomrule
\end{tabular}
\caption{Matrice di tracciabilità Attori/Endpoint}
\label{tab:matrix-endpoints}
\end{table}

% =============================================================================
\chapter{Architettura del Sistema}
% =============================================================================

\section{Visione d'Insieme}

Agent~2 non è un singolo servizio monolitico, ma un sistema modulare in cui
ogni componente ha una responsabilità precisa nella pipeline di validazione:

\begin{itemize}
  \item \textbf{Il Cuore (agent2-api):} Backend FastAPI che espone 12 endpoint,
        inizializza i 4 servizi di analisi all'avvio e orchestra la pipeline
        in 5 step ad ogni chiamata \texttt{/analyze}.

  \item \textbf{L'Orchestratore (n8n):} Gestisce il loop interattivo utente,
        presentando le domande generate dall'agente via form HTML e reinviando
        le risposte all'iterazione successiva tramite webhook.

  \item \textbf{Il Motore AI (Ollama -- host):} Esegue Qwen2.5~14B con accesso
        diretto alla GPU NVIDIA RTX~4070. Non è containerizzato per garantire
        le massime prestazioni; i container lo raggiungono via
        \texttt{host.docker.internal:11434}.

  \item \textbf{Il Consumer Asincrono (agent2-consumer -- futuro):} Modalità
        batch per la pipeline automatica Agent~1 $\rightarrow$ Agent~2
        $\rightarrow$ Agent~3 via Kafka, scalabile a $N$ repliche.

  \item \textbf{Il Bus Eventi (Kafka -- futuro):} Message broker per la
        comunicazione A2A asincrona; topic \texttt{agent1-output} (input) e
        \texttt{agent2-output} (output verso Agent~3).
\end{itemize}

\section{Diagramma Architetturale}
\begin{figure}[H]
\centering
\begin{tikzpicture}[
  node distance=1.4cm and 2.5cm,
  box/.style={draw=agentblue!70, fill=agentblue!8, rounded corners=4pt,
    minimum width=3.8cm, minimum height=1.4cm, align=center,
    font=\small\sffamily},
  mainbox/.style={draw=agentblue, fill=agentblue!15, rounded corners=4pt,
    minimum width=3.8cm, minimum height=1.4cm, align=center,
    font=\small\sffamily\bfseries, line width=1.2pt},
  futurebox/.style={draw=darkgray!50, fill=lightgray, rounded corners=4pt,
    minimum width=3.8cm, minimum height=1.4cm, align=center,
    font=\small\sffamily, dashed},
  arrow/.style={->, >=stealth, thick, color=agentblue!70},
  dashedarrow/.style={->, >=stealth, thick, color=darkgray!60, dashed},
  label/.style={font=\scriptsize\sffamily, color=darkgray}
]
  % Nodes
  \node[box] (agent1) {Agent 1\\[-2pt]\scriptsize Domain Interview};
  \node[box, right=of agent1] (n8n) {n8n\\[-2pt]\scriptsize :5678 (UI)};
  \node[futurebox, below=of agent1] (kafka1) {Kafka: agent1-output\\[-2pt]\scriptsize (futura)};
  \node[mainbox, below=of n8n] (api) {agent2-api\\[-2pt]\scriptsize :8002};
  \node[box, below=of api] (ollama) {Ollama (host)\\[-2pt]\scriptsize qwen2.5:14b -- :11434};
  \node[futurebox, below=of ollama] (kafka2) {Kafka: agent2-output\\[-2pt]\scriptsize (futura)};
  \node[box, right=1.5cm of kafka2] (agent3) {Agent 3\\[-2pt]\scriptsize Spec Generator};

  % Arrows
  \draw[arrow] (agent1) -- node[label, left] {domain\_model JSON} (kafka1);
  \draw[dashedarrow] (kafka1) -- (api);
  \draw[arrow] (n8n) -- node[label, right] {HTTP POST /analyze} (api);
  \draw[arrow] (api) -- (ollama);
  \draw[dashedarrow] (ollama) -- (kafka2);
  \draw[dashedarrow] (kafka2) -- (agent3);
\end{tikzpicture}
\caption{Architettura della pipeline Intelligent Domain Architect con Agent~2 al centro}
\label{fig:arch-overview}
\end{figure}

\section{Topologia Docker}

\subsection{Flusso REST (Modalità Default)}

Il file \texttt{docker-compose.yml} avvia 2 container sulla rete bridge
\texttt{agent2-network} (subnet \texttt{172.28.0.0/16}):

\begin{table}[H]
\centering
\renewcommand{\arraystretch}{1.3}
\begin{tabular}{llp{7cm}}
\toprule
\textbf{Servizio} & \textbf{Porta} & \textbf{Funzione} \\
\midrule
\textbf{agent2-api} & 8002 & Backend FastAPI: 12 endpoint, orchestrazione pipeline 5-step, A2A Protocol v0.3. Build da \texttt{Dockerfile} locale. Limiti: 2 CPU, 2 GB RAM. \\
\textbf{n8n} & 5678 & Orchestratore workflow interattivo. Gestisce il loop ricorsivo (scelta file $\rightarrow$ analisi $\rightarrow$ form domande $\rightarrow$ reinvio risposte $\rightarrow$ ri-analisi fino a VALID). \\
\bottomrule
\end{tabular}
\caption{Servizi core -- docker-compose.yml}
\label{tab:core-services}
\end{table}

\subsection{Flusso Kafka (Predisposto)}

Il file \texttt{docker-compose.kafka.yml} aggiunge 4 container aggiuntivi:

\begin{table}[H]
\centering
\renewcommand{\arraystretch}{1.3}
\begin{tabular}{llp{6.5cm}}
\toprule
\textbf{Servizio} & \textbf{Porta} & \textbf{Funzione} \\
\midrule
\textbf{zookeeper} & 2181 (interna) & Coordinamento broker Kafka (leader election, metadata topic). \\
\textbf{kafka} & 9092/9093 & Message broker A2A. Topic: \texttt{agent1-output} (in), \texttt{agent2-output} (out). 4 partizioni, retention 7 giorni. \\
\textbf{agent2-consumer} & -- & Consumer parallelo da \texttt{agent1-output}. Stessa immagine di agent2-api con entry point dedicato al consumo asincrono. Scalabile via \texttt{CONSUMER\_REPLICAS}. Limiti: 4 CPU, 4 GB RAM. \\
\textbf{kafka-ui} & 8080 & Dashboard web per ispezionare topic, messaggi e consumer groups. \\
\bottomrule
\end{tabular}
\caption{Servizi Kafka -- docker-compose.kafka.yml (predisposto)}
\label{tab:kafka-services}
\end{table}

\section{I Pilastri del Progetto}

\textbf{Separazione delle Responsabilità (SoC):} Ogni fase della pipeline (analisi
semantica, conflict detection, generazione domande, raffinamento) è implementata
in un servizio dedicato. Un fallimento nel modulo LLM non interrompe
l'analisi basata su regole.

\textbf{LLM come Prima Scelta, Euristica come Fallback:} Il sistema prova sempre
l'approccio LLM (Qwen2.5~14B via Ollama) per ottenere la massima qualità.
In caso di timeout, errore o JSON malformato, il fallback garantisce la
continuità del servizio senza intervento umano.

\textbf{Filtraggio Incrementale degli Issue:} Il sistema mantiene memoria delle
decisioni dell'utente nelle iterazioni precedenti (\texttt{\_userGuidedFixes},
\texttt{\_ownershipDecisions}) per non riproporre problemi già risolti.

\textbf{Immagine Docker Unica, Due Entry Point:} \texttt{agent2-api} e
\texttt{agent2-consumer} usano la stessa immagine \texttt{agent2-consistency-analyzer:latest},
garantendo parità di comportamento tra modalità REST e Kafka.

\section{Architettura dei Servizi Interni}

Il backend di Agent~2 è composto da quattro servizi indipendenti, inizializzati
all'avvio dell'applicazione e iniettati come dipendenze agli endpoint.
Ogni servizio è responsabile di una singola fase della pipeline.

\subsection{SemanticAnalyzer}

Il \texttt{SemanticAnalyzer} è il primo stadio della pipeline e opera interamente
senza LLM, basandosi su calcolo vettoriale puro.

\textbf{Inizializzazione:} Al primo avvio, scarica e carica il modello
\texttt{all-MiniLM-L6-v2} da HuggingFace nella directory \texttt{./models}.
Questo avviene una sola volta; i successivi avvii usano la cache locale.
Il modello produce embedding a \textbf{384 dimensioni}.

\textbf{Algoritmo di rilevamento ENTITY\_OVERLAP:}
\begin{enumerate}
  \item Estrae tutte le entità da ogni bounded context del \texttt{domainMap}.
  \item Calcola l'embedding di ogni nome di entità.
  \item Calcola la \textit{cosine similarity} tra ogni coppia di entità
        appartenenti a bounded context diversi.
  \item Se la similarità supera \texttt{ENTITY\_OVERLAP\_THRESHOLD} (default: 0.85),
        genera un issue di tipo \texttt{ENTITY\_OVERLAP} con i due contesti
        coinvolti e il punteggio di similarità.
\end{enumerate}

\textbf{Algoritmo di rilevamento SEMANTIC\_AMBIGUITY:}
\begin{enumerate}
  \item All'interno dello stesso bounded context, confronta entità con
        definizioni presenti nell'\texttt{ubiquitousLanguage}.
  \item Se due definizioni hanno similarità superiore a
        \texttt{SEMANTIC\_SIMILARITY\_THRESHOLD} (default: 0.75) pur riferendosi
        a concetti che dovrebbero essere distinti, genera un issue
        \texttt{SEMANTIC\_AMBIGUITY}.
\end{enumerate}

\begin{technicalbox}{Dettaglio Tecnico: Filtraggio Pre-Analisi}
Prima di calcolare gli embedding, il servizio esclude automaticamente:
\begin{itemize}
  \item Entità con flag \texttt{\_referenceOnly: true} (già risolte dallo Step~0)
  \item Issue i cui \texttt{affected\_elements} compaiono già in
        \texttt{\_userGuidedFixes} o \texttt{\_ownershipDecisions} nel modello
\end{itemize}
Questo filtraggio garantisce che le iterazioni successive analizzino solo
i problemi ancora aperti, evitando di riproporre questioni già risolte
dall'utente nelle iterazioni precedenti.
\end{technicalbox}

\subsection{ConflictDetector}

Il \texttt{ConflictDetector} implementa un approccio ibrido: prima applica
22 regole DDD deterministiche, poi usa l'LLM per approfondire i casi ambigui.

\textbf{Fase 1 -- Regole DDD:} Ogni regola della \texttt{validation\_checklist.json}
viene verificata strutturalmente sul modello. Esempio: la regola EA-001
(``ogni evento deve avere un solo emitter'') scorre tutti gli eventi e verifica
che il campo \texttt{emitter} sia unico; se un evento appare in più domini
come emitter, genera \texttt{DUPLICATE\_EVENT}.

\textbf{Fase 2 -- Analisi LLM:} Il servizio costruisce un prompt che include
le 22 regole DDD come contesto (da \texttt{ddd\_rules.md}) e la rappresentazione
JSON del modello, chiedendo a Qwen2.5~14B di identificare conflitti non
catturati dalle regole formali. La risposta viene parsata e convertita in
\texttt{ConflictIssue}.

\begin{technicalbox}{Dettaglio Tecnico: Categorie Knowledge Base}
Le 7 categorie di regole nella \texttt{validation\_checklist.json}:
\begin{center}
\begin{tabular}{lll}
\toprule
Codice & Categoria & N. regole \\
\midrule
DC  & Domain Classification        & 3 \\
BC  & Bounded Context              & 3 \\
EVO & Entity vs Value Object       & 3 \\
EO  & Entity Ownership             & 3 \\
RC  & Requirement Consistency      & 3 \\
EA  & Event Architecture           & 4 \\
CP  & Communication Patterns       & 3 \\
\midrule
    & \textbf{Totale}              & \textbf{22} \\
\bottomrule
\end{tabular}
\end{center}
\end{technicalbox}

\subsection{QuestionGenerator}

Il \texttt{QuestionGenerator} traduce gli issue tecnici rilevati in domande
comprensibili per il business stakeholder, con risposte suggerite azionabili.

\textbf{Approccio LLM-first (v2.0):} Tutti gli issue vengono inviati in una
singola chiamata a Qwen2.5~14B, che genera domande specifiche usando i nomi reali
delle entità del modello. Il numero massimo di domande è configurabile tramite
\texttt{MAX\_FOLLOW\_UP\_QUESTIONS} (default: 5).

\textbf{Struttura di ogni domanda generata:}
\begin{itemize}
  \item \texttt{question\_id}: Identificativo univoco (es. \texttt{FUQ-001})
  \item \texttt{question}: Testo della domanda con nomi reali di entità e contesti
  \item \texttt{severity}: Livello di urgenza (\texttt{HIGH}/\texttt{MEDIUM}/\texttt{LOW})
  \item \texttt{related\_issue\_type}: Tipo di issue DDD correlato
  \item \texttt{affected\_elements}: Lista di entità/contesti coinvolti
  \item \texttt{suggested\_answers}: 2--3 risposte suggerite azionabili
\end{itemize}

\textbf{Fallback Template:} Se il LLM non è disponibile, il servizio usa i
\texttt{question\_template} presenti in ogni regola della
\texttt{validation\_checklist.json}, sostituendo i placeholder con i nomi
reali delle entità coinvolte.

\subsection{ModelRefiner}
\label{sec:model-refiner-detail}

Il \texttt{ModelRefiner} è il servizio più complesso: gestisce sia gli auto-fix
deterministici che l'interpretazione LLM delle risposte utente.

\textbf{Auto-fix deterministici:} Per \texttt{NAMING\_VIOLATION}, il refiner
converte automaticamente i nomi degli eventi dal formato imperativo al passato
(es. \texttt{CreateOrder} $\rightarrow$ \texttt{OrderCreated}) senza richiedere
conferma utente.

\textbf{Interpretazione delle risposte (catena a 3 livelli):} Quando l'utente
risponde alle domande di follow-up, ogni risposta attraversa una catena di
interpretazione progressiva:
\begin{enumerate}
  \item \textbf{LLM (prioritario):} La risposta viene inviata al modello LLM
        con il contesto dell'issue correlato. Il modello restituisce una delle
        6 azioni strutturate (\texttt{rename\_event}, \texttt{set\_owner}, ecc.)
        con i parametri specifici. Se il LLM restituisce \texttt{no\_change},
        il sistema riprova fino a 3 tentativi.
  \item \textbf{Rule-based fallback:} Se il LLM non è disponibile o fallisce,
        il sistema tenta un matching basato su regole: prima confronta la risposta
        con le \texttt{suggested\_answers} generate (fuzzy matching), poi applica
        pattern matching specifici per tipo di issue (es. regex per estrarre nomi
        di contesto da risposte \texttt{ENTITY\_OVERLAP}, o accettazione diretta
        per \texttt{SEMANTIC\_AMBIGUITY}).
  \item \textbf{Catch-all:} Se entrambi i livelli precedenti falliscono e l'utente
        ha fornito una risposta non vuota, il sistema registra una risoluzione
        generica (\texttt{resolve\_conflict}) per evitare loop infiniti di
        re-ask della stessa domanda.
\end{enumerate}
Ogni interpretazione viene tracciata con il campo \texttt{interpretation\_source}
(\texttt{llm}, \texttt{rules} o \texttt{catch-all}) e lo stato \texttt{applied}
(booleano onesto: \texttt{true} solo se la modifica strutturale è stata effettivamente
applicata al modello).

\textbf{Annotazioni nel modello:} Le modifiche vengono tracciate nei metadati
del modello:
\begin{itemize}
  \item \texttt{\_userGuidedFixes}: storia delle risoluzioni manuali per ownership
  \item \texttt{\_ownershipDecisions}: assegnazioni di ownership approvate dall'utente
  \item \texttt{\_conflictResolutions}: conflitti esplicitamente risolti
  \item \texttt{\_validationAnnotations}: annotazioni dell'ultima iterazione
        (rimosse all'inizio della successiva dallo Step~0)
\end{itemize}

\section{Perché Queste Tecnologie?}
\begin{technicalbox}{Dettaglio Tecnico: Stack Tecnologico}
\textbf{FastAPI:} Framework REST asincrono con generazione automatica di
documentazione OpenAPI/Swagger e validazione degli schemi. Ideale per gestire
payload JSON complessi e variabili.

\textbf{Qwen2.5~14B (Ollama):} $\sim$9 GB in formato Q4\_K\_M, inferenza in VRAM
su RTX~4070. Eccellente nella generazione di JSON strutturato e nel ragionamento
multi-step, con prestazioni superiori a modelli da 7B su task di analisi e classificazione.

\textbf{all-MiniLM-L6-v2:} Modello sentence-transformers (23M parametri) con
embedding a 384 dimensioni, ottimizzato per rilevare sovrapposizioni semantiche
tra entità DDD.

\textbf{Kafka (Confluent 7.5.0):} Standard de facto per pipeline event-driven
A2A. 4 partizioni per topic garantiscono parallelismo; retention 7 giorni
permette replay in caso di fallimenti downstream.

\textbf{n8n:} Workflow engine low-code per gestire il loop interattivo tramite
webhook e form HTML, senza necessità di un frontend dedicato.
\end{technicalbox}

% =============================================================================
\chapter{Flussi di Funzionamento}
% =============================================================================

\section{Pipeline di Analisi -- Visione d'Insieme}

Ogni chiamata a \texttt{POST /analyze} esegue fino a 5 step in sequenza.
Se sono presenti risposte utente da un'iterazione precedente
(\texttt{previous\_answers}), viene eseguito prima lo \textbf{Step 0}.

\subsection{Step 0 -- Pre-applicazione Risposte (sole iterazioni successive)}

Quando \texttt{/analyze} viene chiamato con \texttt{previous\_answers}
(dalla seconda iterazione in poi nel loop n8n), le risposte dell'utente
vengono \textbf{applicate al domain model PRIMA} di eseguire l'analisi.

\begin{technicalbox}{Dettaglio Tecnico: Perché Prima dell'Analisi}
Applicare le risposte \textit{dopo} l'analisi (come nella v1.x) comportava
che i problemi già risolti dall'utente venissero rilevati nuovamente,
rendendo il loop inutile. Con lo Step 0, il modello parte da uno stato
aggiornato e l'analisi rileva solo i problemi ancora aperti.

\textbf{Operazioni di Step 0:}
\begin{itemize}
  \item Recupera \texttt{\_followUpQuestions} memorizzate nel modello
        dall'iterazione precedente.
  \item Invoca il servizio ModelRefiner per
        interpretare ogni risposta via LLM e applicare le modifiche concrete
        (\texttt{rename\_event}, \texttt{set\_owner}, \texttt{reclassify\_domain},
        \texttt{resolve\_conflict}, \texttt{add\_mapping}, \texttt{update\_pattern}).
  \item Rimuove i metadati dell'iterazione precedente
        (\texttt{\_validationAnnotations}, \texttt{\_followUpQuestions})
        per partire puliti; conserva \texttt{\_userGuidedFixes} e
        \texttt{\_ownershipDecisions} per il filtraggio degli issue futuri.
\end{itemize}
\end{technicalbox}

\subsection{Step 1 -- Analisi Semantica (Embeddings)}

\textbf{Servizio:} SemanticAnalyzer

Usa il modello \texttt{all-MiniLM-L6-v2} (sentence-transformers) per calcolare
la similarità coseno tra entità di bounded context diversi.

\textbf{Rileva:}
\begin{itemize}
  \item \texttt{ENTITY\_OVERLAP} -- stessa entità in più bounded context
        senza differenziazione (\textit{threshold}: 0.85)
  \item \texttt{SEMANTIC\_AMBIGUITY} -- termini con significati diversi
        nello stesso contesto (\textit{threshold}: 0.75)
\end{itemize}

\begin{technicalbox}{Dettaglio Tecnico: Filtraggio Intelligente}
Le entità marcate \texttt{\_referenceOnly} (già assegnate a un owner dallo Step 0)
vengono escluse dall'analisi. Gli issue i cui \texttt{affected\_elements}
sono già presenti in \texttt{\_userGuidedFixes} o \texttt{\_ownershipDecisions}
vengono scartati automaticamente, evitando di riproporre problemi già risolti.
\end{technicalbox}

\subsection{Step 2 -- Conflict Detection (Regole + LLM)}

\textbf{Servizio:} ConflictDetector

Applica 22 regole DDD dalla knowledge base (\texttt{ddd\_rules.md}) suddivise
in 7 categorie, con analisi LLM aggiuntiva per i casi complessi.

\textbf{Rileva:}
\begin{itemize}
  \item \texttt{REQUIREMENT\_CONFLICT} -- requisiti che si contraddicono
  \item \texttt{DUPLICATE\_EVENT} -- stesso evento emesso da più domini
  \item \texttt{NAMING\_VIOLATION} -- eventi non al tempo passato
        (es.\ \texttt{CreateOrder} anziché \texttt{OrderCreated})
  \item \texttt{INCOMPATIBLE\_PATTERN} -- pattern sincrono richiesto su
        canale solo asincrono
  \item \texttt{MISCLASSIFIED\_DOMAIN} -- dominio Core classificato
        erroneamente come Generic/Supporting
\end{itemize}

\begin{technicalbox}{Dettaglio Tecnico: Parametri LLM per Servizio}
\begin{center}
\begin{tabular}{lccc}
\toprule
Servizio & Temperature & Max tokens & Timeout \\
\midrule
Conflict Detector & 0.3 & 2048 & 120s \\
Question Generator & 0.4 & 4096 & 120s \\
Model Refiner & 0.2 & 2048 & 120s \\
\bottomrule
\end{tabular}
\end{center}
Temperature bassa nel refiner (0.2) per massimo determinismo nelle
modifiche al modello; più alta nel generator (0.4) per domande contestualmente
variate.
\end{technicalbox}

\subsection{Step 3 -- Generazione Domande (LLM-first)}

\textbf{Servizio:} QuestionGenerator

Approccio \textbf{LLM-first} (v2.0): quando \texttt{use\_llm=true},
\textbf{tutte} le domande vengono generate da Qwen2.5~14B in una singola chiamata.
Il modello riceve tutti gli issue rilevati e produce domande specifiche con
nomi reali delle entità e risposte suggerite azionabili.

\textbf{Esempio di output LLM:} Il sistema genera una domanda contestuale
come \textit{``Come gestire l'ownership di Product tra CatalogContext
e OrderContext?''}, accompagnata da 2--3 risposte suggerite azionabili
(es.\ assegnare ownership a un contesto specifico, rinominare le entità
per distinguere i ruoli, o consolidare in un unico bounded context).

\textbf{Fallback:} se il LLM fallisce (timeout, JSON malformato), usa i
\texttt{question\_template} dalla \texttt{validation\_checklist.json}.

\subsection{Step 4 -- Raffinamento Modello}

\textbf{Servizio:} ModelRefiner

Auto-fix per problemi semplici (naming violations $\rightarrow$ conversione al passato).
Nelle iterazioni successive, interpreta le risposte utente e applica modifiche
strutturali al modello:

\begin{table}[H]
\centering
\renewcommand{\arraystretch}{1.5}
\begin{tabular}{c @{\hspace{1cm}} p{10cm}}
\toprule
\textbf{Azione LLM} & \textbf{Modifica applicata al modello} \\
\midrule
\texttt{rename\_event}      & Rinomina l'evento nel modello \\
\texttt{set\_owner}         & Assegna ownership, sostituisce l'entità con riferimento nei contesti non-owner \\
\texttt{reclassify\_domain} & Sposta il dominio tra core/supporting/generic \\
\texttt{resolve\_conflict}  & Registra la risoluzione del conflitto \\
\texttt{add\_mapping}       & Aggiunge un mapping anti-corruption layer \\
\texttt{update\_pattern}    & Aggiorna il pattern di integrazione \\
\bottomrule
\end{tabular}
\caption{Azioni di raffinamento interpretate dall'LLM}
\label{tab:refiner-actions}
\end{table}

\section{Stima Adattiva delle Iterazioni}

Dopo i 4 step, Agent~2 stima il numero di iterazioni necessarie per raggiungere
lo stato \texttt{VALID}, restituito nel campo \texttt{suggested\_max\_iterations}.

\begin{technicalbox}{Dettaglio Tecnico: Strategia a Due Livelli}
\textbf{LLM (prioritario):} Qwen2.5 riceve un riepilogo dei problemi (quantità,
severità, tipologia) e restituisce un intero tra 1 e 10 con
\texttt{temperature: 0.1} per massima determinismo.

\textbf{Euristica (fallback):}
\begin{center}
\begin{tabular}{ccc}
\toprule
Problemi manuali & Issue HIGH & Iterazioni stimate \\
\midrule
0     & -- & 1 \\
1     & -- & 1 \\
2--3  & -- & 2 \\
4--5  & -- & 3 \\
6+    & $\geq 3$ & min(N, 7) \\
\bottomrule
\end{tabular}
\end{center}

Il workflow n8n aggiorna \texttt{max\_iterations} con l'ultimo valore stimato
ad ogni iterazione (soft cap adattivo): se il contatore raggiunge il limite,
\textit{Prepara Dati} lo incrementa automaticamente di 1 e il loop prosegue.
\end{technicalbox}
\section{Workflow n8n -- Loop Interattivo}
Il file \texttt{workflow\_complete\_loop.json} implementa un workflow a 15 nodi
che gestisce l'intero ciclo di vita dell'analisi:
\begin{figure}[H]
\centering
\tikzset{
  block/.style    = {rectangle, rounded corners, draw=blue!60, fill=blue!10,
                     minimum width=3.2cm, minimum height=0.7cm, align=center, font=\small},
  decision/.style = {diamond, draw=orange!70, fill=orange!15, aspect=2.5,
                     align=center, font=\small},
  arrow/.style    = {-{Stealth}, thick}
}
\begin{tikzpicture}[x=1cm, y=1cm]

  % colonna centrale (x=0)
  \node[block]    at (0,  0)    (start)      {Browser GET\\/webhook/agent2-start};
  \node[block]    at (0, -2)    (scelta)     {Pagina Scelta};
  \node[decision] at (0, -4)    (filelocale) {File Locale?};
  \node[block]    at (0, -6.5)  (prepara)    {Prepara Dati};
  \node[decision] at (0, -9)    (maxiter)    {Max Iter?};
  \node[decision] at (0, -14)   (valid)      {Modello Valido?};

  % colonna destra (x=5)
  \node[block]    at (5, -4)    (lista)      {Lista File\\+ Selezione};
  \node[block]    at (5, -9)    (analyze)    {POST /analyze};
  \node[block]    at (5, -11)   (combina)    {Combina Risultati};
  \node[block]    at (5, -12.6) (form)       {Form Domande HTML};

  % colonna destra estrema (x=9)
  \node[block]    at (9, -15.5) (successo)   {Pagina Successo};

  % frecce colonna centrale
  \draw[arrow] (start)  -- (scelta);
  \draw[arrow] (scelta) -- (filelocale);
  \draw[arrow] (filelocale) -- node[right]{\small No} (prepara);
  \draw[arrow] (prepara) -- (maxiter);

  % max iter: continua con analyze (soft cap)
  \draw[arrow] (maxiter.east) -- node[above]{\small soft cap} (analyze.west);

  % max iter: cap reached -> vai a valid (dritto, più pulito)
  \draw[arrow] (maxiter.south) -- node[right]{\small cap reached} (valid.north);

  % ramo file locale = sì
  \draw[arrow] (filelocale.east) -- node[above]{\small Sì} (lista.west);
  % lista -> prepara: una sola freccia con angolo retto
  \draw[arrow] (lista.south) -- ++(0,-2) coordinate (ls_turn) |-
  (prepara.east);

  % analisi e form
  \draw[arrow] (analyze) -- (combina);
  \draw[arrow] (combina) -- (form);

  % form -> valid (post submit) entra da destra in valid
  \draw[arrow] (form.east) -- ++(1.2,0)
              |- node[right, pos=0.25]{\small POST /submit} (valid.east);

  % valid -> successo (sì) parte dal rombo, scende fino alla y di successo, poi va a destra
  \draw[arrow] (valid.south) -- ++(0,-0.3)
              |-node[above=2pt, pos=0.7]{\small Sì} (successo.west);

  % loop no: valid -> prepara (lato sinistro)
  \draw[arrow] (valid.west) -- ++(-2,0)
              |- node[left, pos=0.25]{\small No} (prepara.west);

\end{tikzpicture}
\caption{Workflow n8n -- loop interattivo completo (15 nodi)}
\label{fig:n8n-loop}
\end{figure}



\begin{table}[H]
\centering
\renewcommand{\arraystretch}{1.3}
\begin{tabularx}{\textwidth}{l l X}
\toprule
\textbf{Nodo} & \textbf{Tipo} & \textbf{Funzione} \\
\midrule
START               & Webhook GET      & Pagina iniziale nel browser \\
SUBMIT              & Webhook POST     & Riceve tutte le form submissions \\
Detect Request Type & Code             & Classifica la richiesta (file/answers/initial) \\
Route               & Switch           & Smista il flusso verso GET start, POST submit o risposta file \\
Pagina Scelta       & Respond Webhook  & Serve l'HTML con i bottoni ``File locale'' / ``Incolla JSON'' \\
Lista File          & HTTP Request     & Legge i file \texttt{.json} presenti in \texttt{input\_agent/} \\
Selezione File      & Respond Webhook  & Serve l'HTML con la lista dei file selezionabili \\
Carica File         & HTTP Request     & Legge dal disco il file JSON selezionato dall'utente \\
Prepara Dati        & Code             & Normalizza iteration/sessione, applica soft cap \\
Analizza Modello    & HTTP Request     & \texttt{POST http://agent2-api:8002/analyze} \\
Combina Risultati   & Code             & Pre-renderizza HTML domande, aggiorna \texttt{max\_iterations} \\
Modello Valido?     & IF               & Controlla \texttt{status == VALID}: sì $\to$ successo, no $\to$ form \\
Form Domande        & Respond Webhook  & Mostra domande con radio buttons e testo libero \\
Pagina Successo     & Respond Webhook  & Risultato finale quando \texttt{status=VALID} \\
A2A Wait            & Respond Webhook  & Gestisce il ramo agente-a-agente (flusso non interattivo) \\
\bottomrule
\end{tabularx}
\caption{Nodi del workflow n8n (15 totali)}
\label{tab:n8n-nodes}
\end{table}

\section{Esempio Completo: Richiesta \texttt{POST /analyze}}

\subsection{Struttura della Richiesta}

Il campo obbligatorio è \texttt{domain\_model}; \texttt{use\_llm} e
\texttt{previous\_answers} sono opzionali (default: \texttt{true} e \texttt{[]}).

La richiesta \texttt{POST /analyze} riceve un domain model strutturato
contenente: metadati del modello (identificativo, versione, autore),
contesto di business, mappa dei domini (core, supporting, generic) con
le relative entità, value object e requisiti funzionali, modello
event-driven con gli eventi di dominio, e le integrazioni tra bounded
context con i relativi pattern di comunicazione.

\textbf{Esempio di scenario con issue intenzionali:} Si consideri un
modello e-commerce con due domini core (\texttt{OrderManagement} e
\texttt{CatalogManagement}), ciascuno contenente un'entità
\texttt{Product} con significati diversi, un evento denominato
\texttt{CreateOrder} (in forma imperativa anziché al passato) e
requisiti funzionali contraddittori tra i due domini.

\textbf{Issue attesi in questo scenario:}
\begin{itemize}
  \item \texttt{ENTITY\_OVERLAP} (HIGH): \texttt{Product} appare sia in
        \texttt{OrderContext} che in \texttt{CatalogContext} senza differenziazione.
  \item \texttt{REQUIREMENT\_CONFLICT} (HIGH): FR-001 (\textit{ordine immutabile})
        e FR-002 (\textit{prodotto modificabile in qualsiasi momento}) si contraddicono
        poiché l'ordine contiene un riferimento al prodotto.
  \item \texttt{NAMING\_VIOLATION} (MEDIUM): \texttt{CreateOrder} non è al tempo
        passato (dovrebbe essere \texttt{OrderCreated}).
\end{itemize}

\subsection{Struttura della Risposta}

La risposta di Agent~2 contiene i seguenti elementi principali:

\begin{itemize}
  \item \textbf{Task ID:} Identificativo univoco UUID della sessione di analisi.
  \item \textbf{Status:} Stato del modello (\texttt{VALID} oppure \texttt{ISSUES\_FOUND}).
  \item \textbf{Summary:} Riepilogo numerico con totale issue, distribuzione per severità
        (HIGH/MEDIUM/LOW), numero di domini e entità analizzati.
  \item \textbf{Issues:} Lista dettagliata di ogni problema rilevato, con tipologia,
        severità, descrizione testuale, elementi coinvolti e suggerimento di fix.
        Gli issue auto-corretti (es.\ naming violations) sono marcati come tali.
  \item \textbf{Follow-up Questions:} Domande contestuali generate dal LLM, ciascuna
        con identificativo univoco, testo, severità, tipo di issue correlato e
        2--3 risposte suggerite azionabili.
  \item \textbf{Refined Model:} Versione aggiornata del domain model con le correzioni
        automatiche già applicate (es.\ rinomina eventi al tempo passato).
  \item \textbf{Suggested Max Iterations:} Stima adattiva del numero di iterazioni
        necessarie per raggiungere lo stato \texttt{VALID}.
\end{itemize}

\section{Ciclo di Vita della Richiesta REST}

Ogni chiamata HTTP ad Agent~2 segue una pipeline standard:

\begin{technicalbox}{Dettaglio Tecnico: Gestione della Richiesta}
\begin{enumerate}
  \item \textbf{Ricezione:} Il framework deserializza il payload JSON
        con validazione automatica dello schema.
  \item \textbf{Task ID:} Viene generato un UUID univoco come identificativo
        di sessione per il tracciamento del risultato.
  \item \textbf{Pipeline:} Esecuzione sequenziale dei 5 step di analisi.
  \item \textbf{Persistenza:} Il risultato viene salvato in
        \texttt{output\_agent/\{filename\}/\{session\_id\}/iteration\_\{N\}.json}.
  \item \textbf{Risposta:} Restituzione della \texttt{ValidationResponse}
        con status, summary, issue, domande e modello raffinato.
\end{enumerate}
\end{technicalbox}
\section{Gestione della Resilienza}

\begin{technicalbox}{Dettaglio Tecnico: Error Handling e Fallback AI}
\textbf{Exception Handling:} Le eccezioni vengono catturate a livello
di endpoint e restituite come payload JSON strutturati con codice HTTP 500.

\textbf{Fallback AI:} In caso di timeout o errore di Ollama:
\begin{itemize}
  \item \textit{Conflict Detector} -- usa solo le 22 regole DDD (nessun LLM)
  \item \textit{Question Generator} -- usa i \texttt{question\_template}
        dalla \texttt{validation\_checklist.json}
  \item \textit{Model Refiner} -- utilizza la catena di interpretazione a 3 livelli
        (Sezione~\ref{sec:model-refiner-detail}): se il LLM non è disponibile,
        il fallback rule-based gestisce le risposte tramite matching delle
        \texttt{suggested\_answers} e pattern specifici per tipo di issue
        (incluso \texttt{SEMANTIC\_AMBIGUITY}); come ultima risorsa, il catch-all
        registra la decisione utente per evitare loop infiniti
  \item \textit{Iteration Estimator} -- usa la formula euristica
\end{itemize}
In tutti i casi il servizio risponde, garantendo la continuità del flusso.
\end{technicalbox}

% =============================================================================
\chapter{Architettura e Governance del Dato}
% =============================================================================

\section{Modello Dati del Domain Model}

Il dato centrale elaborato da Agent~2 è il \texttt{domain\_model} JSON,
struttura gerarchica composta da:

\begin{itemize}
  \item \textbf{metadata} -- identificativo del modello, versione, autore (Agent~1),
        timestamp di creazione.
  \item \textbf{businessContext} -- obiettivi di business e perimetro del dominio.
  \item \textbf{domainMap} -- mappa dei domini suddivisa in
        \texttt{coreDomains}, \texttt{supportingDomains}, \texttt{genericDomains}.
        Ogni dominio contiene \texttt{entities}, \texttt{valueObjects},
        \texttt{functionalRequirements} e \texttt{boundedContext} con
        \texttt{ubiquitousLanguage}.
  \item \textbf{eventDrivenModel} -- lista degli eventi di dominio con emitter,
        consumer, payload e pattern di comunicazione.
  \item \textbf{contextIntegrations} -- relazioni tra bounded context con
        tipo di integrazione e pattern comunicazione.
\end{itemize}

\begin{technicalbox}{Dettaglio Tecnico: Schema di Relazione tra Entità DDD}
\textbf{One-to-Many:} Un dominio possiede N entità; un evento ha un emitter
e N consumer. \\
\textbf{Cross-Context Reference:} Dopo lo Step 0, le entità con ownership
assegnata vengono sostituite con riferimenti (\texttt{\_referenceOnly: true})
nei contesti non-owner, mantenendo la struttura originale come storia. \\
\textbf{Metadata di Tracciabilità:} I campi \texttt{\_userGuidedFixes},
\texttt{\_ownershipDecisions} e \texttt{\_conflictResolutions} vengono
accumulati iterazione per iterazione per tracciare la storia delle decisioni.
\end{technicalbox}

\section{Schema Completo del Domain Model}

Il domain model è il contratto principale tra Agent~1, Agent~2 e Agent~3.
Lo schema è composto dalle seguenti sezioni principali:

\begin{table}[H]
\centering
\renewcommand{\arraystretch}{1.3}
\begin{tabular}{lp{10cm}}
\toprule
\textbf{Sezione} & \textbf{Contenuto} \\
\midrule
\texttt{metadata}             & Identificativo del modello, versione, autore (Agent~1), timestamp di creazione. \\
\texttt{businessContext}      & Descrizione del business, obiettivi e perimetro del dominio. \\
\texttt{domainMap}            & Mappa dei domini suddivisa in \texttt{coreDomains}, \texttt{supportingDomains} e \texttt{genericDomains}. \\
\texttt{eventDrivenModel}     & Lista degli eventi di dominio con emitter, consumer, payload e pattern di comunicazione. Gli eventi rinominati da Agent~2 sono marcati come auto-corretti. \\
\texttt{contextIntegrations}  & Relazioni tra bounded context con tipo di integrazione (sync/async) e pattern (REST, Kafka, gRPC, GraphQL). \\
\bottomrule
\end{tabular}
\caption{Struttura principale del Domain Model}
\label{tab:domain-model-schema}
\end{table}

\textbf{Metadati aggiunti iterativamente da Agent~2:}
\begin{itemize}
  \item \texttt{\_userGuidedFixes}: storia delle risoluzioni manuali con ID domanda, azione applicata e timestamp.
  \item \texttt{\_ownershipDecisions}: mappa delle assegnazioni di ownership approvate dall'utente.
  \item \texttt{\_conflictResolutions}: conflitti risolti con identificativo, descrizione e modalità (utente o automatica).
  \item \texttt{\_validationAnnotations}: annotazioni temporanee dell'iterazione corrente, rimosse dallo Step~0 successivo.
  \item \texttt{\_followUpQuestions}: domande generate nell'iterazione corrente per il workflow n8n.
\end{itemize}

\textbf{DomainDefinition} -- ogni dominio nell'array contiene:
\begin{itemize}
  \item \textbf{Nome e classificazione:} identificativo del dominio e tipo (core, supporting, generic).
  \item \textbf{Entità:} lista di entità con nome e attributi. Dopo l'assegnazione di ownership, le entità nei contesti non-owner vengono marcate come riferimenti.
  \item \textbf{Value Object e Aggregati:} strutture DDD con i relativi attributi e radici.
  \item \textbf{Requisiti:} funzionali (es.\ FR-001) e non funzionali (es.\ NFR-001).
  \item \textbf{Bounded Context:} nome del contesto, ubiquitous language (dizionario termine-definizione) e dipendenze esterne.
\end{itemize}

\section{Entità di Supporto e Tracciabilità}

\begin{itemize}
  \item \textbf{ValidationResponse:} Il risultato completo di ogni analisi,
        comprendente summary numerico, lista degli issue per categoria,
        follow-up questions e modello raffinato. Persistito come JSON
        in \texttt{output\_agent/}.

  \item \textbf{Task (A2A):} Entità in-memory che traccia lo stato
        (\texttt{SUBMITTED} $\rightarrow$ \texttt{WORKING} $\rightarrow$
        \texttt{COMPLETED}/\texttt{FAILED}/\texttt{CANCELLED}) di ogni
        richiesta ricevuta via protocollo A2A.

  \item \textbf{SemanticIssue / ConflictIssue:} Entità intermedie generate
        rispettivamente dal SemanticAnalyzer e
        dal ConflictDetector, con tipologia, severità ed
        elementi affected. Serializzate in formato JSON per la risposta.

  \item \textbf{FollowUpQuestion:} Struttura dati che rappresenta una singola
        domanda contestuale con ID univoco, testo, severità, tipo di issue
        correlato e lista di risposte suggerite.
\end{itemize}

\section{Knowledge Base}

La knowledge base è il fondamento delle regole DDD applicate nello Step~2
(la suddivisione in 7 categorie è riportata nella Tabella del Capitolo~6):

\begin{itemize}
  \item \textbf{\texttt{ddd\_rules.md} (21 KB):} Guida DDD strutturata in
        8 parti: classificazione domini, bounded context, entity vs value object,
        regole di coerenza tra domini, event-driven architecture, checklist di
        validazione, tipologie di problemi da rilevare e template per l'output.
        Usata dall'LLM come contesto nel prompt di conflict detection.

  \item \textbf{\texttt{validation\_checklist.json}:} 22 regole in formato
        strutturato JSON, organizzate in 7 categorie (DC, BC, EVO, EO, RC, EA, CP).
        Ogni regola include un \texttt{question\_template} per il fallback del
        question generator.
\end{itemize}
\section{Tipologie di Issue Rilevate}

\begin{table}[H]
\centering
\renewcommand{\arraystretch}{1.3}
\begin{tabular}{llp{7cm}}
\toprule
\textbf{Tipo} & \textbf{Severità} & \textbf{Descrizione} \\
\midrule
\texttt{ENTITY\_OVERLAP}       & HIGH   & Stessa entità in più bounded context senza differenziazione. \\
\texttt{SEMANTIC\_AMBIGUITY}   & MEDIUM & Termine con significati diversi nel contesto. \\
\texttt{REQUIREMENT\_CONFLICT} & HIGH   & Requisiti funzionali che si contraddicono. \\
\texttt{DUPLICATE\_EVENT}      & HIGH   & Stesso evento emesso da più domini. \\
\texttt{NAMING\_VIOLATION}     & LOW--MEDIUM & Evento non denominato al tempo passato. \\
\texttt{INCOMPATIBLE\_PATTERN} & HIGH   & Pattern comunicazione incompatibile con i requisiti. \\
\texttt{MISCLASSIFIED\_DOMAIN} & MEDIUM & Classificazione dominio errata (es.\ Auth come Core). \\
\bottomrule
\end{tabular}
\caption{Tipologie di issue rilevate da Agent~2}
\label{tab:issue-types}
\end{table}

\section{Persistenza e Integrità}

\begin{technicalbox}{Dettaglio Tecnico: Strategia di Output}
\textbf{File-based Output:} Ogni analisi viene persistita in
\texttt{output\_agent/\{filename\}/\{session\_id\}/iteration\_\{N\}.json},
dove \texttt{filename} è il nome del file di input (o del modello),
\texttt{session\_id} è l'identificativo di sessione (o timestamp UTC)
e \texttt{N} è il numero di iterazione a due cifre.
La directory viene creata automaticamente se assente.

\textbf{In-Memory Task Store:} I task A2A sono conservati in
una struttura dati in memoria per la durata del processo. Per ambienti
di produzione è prevista la sostituzione con Redis o un database esterno.

\textbf{Read-Only Knowledge Base:} Il volume \texttt{./knowledge\_base}
è montato in sola lettura (\texttt{:ro}) nel container, prevenendo
modifiche accidentali alle regole DDD durante l'esecuzione.
\end{technicalbox}

% =============================================================================
\chapter{Sicurezza e Governance}
% =============================================================================

\section{Modello di Accesso}

Agent~2 opera in un contesto di pipeline interna: gli endpoint non sono
esposti direttamente a utenti finali anonimi ma accessibili da:
\begin{itemize}
  \item n8n (sulla rete interna Docker \texttt{agent2-network})
  \item Agent~1/Agent~3 via A2A Protocol con opzionale JWT Bearer
  \item Team DevOps per debug e configurazione
\end{itemize}

\begin{technicalbox}{Dettaglio Tecnico: Meccanismi di Sicurezza}
\textbf{Container non-root:} Il Dockerfile crea l'utente \texttt{agent2user}
(UID 1000) e avvia l'applicazione con privilegi ridotti.

\textbf{Network Isolation:} La rete bridge \texttt{agent2-network}
(subnet \texttt{172.28.0.0/16}) isola i container dai servizi esterni.
Solo le porte 8002 (agent2-api) e 5678 (n8n) sono esposte all'host.

\textbf{A2A Authentication:} L'Agent Card definisce due schemi di sicurezza:
JWT Bearer Token e API Key (header \texttt{X-API-Key}), entrambi opzionali
nella configurazione attuale ma predisposti per ambienti più restrittivi.

\textbf{Read-Only Volumes:} \texttt{./knowledge\_base} e \texttt{./input\_agent}
sono montati in sola lettura, limitando la superficie di attacco.

\textbf{CORS:} Abilitato per tutti i metodi e header (\texttt{allow\_origins=["*"]})
nella configurazione corrente, adeguata per la pipeline interna.
\end{technicalbox}

\section{Governance dell'Analisi}

La governance assicura che solo modelli validati procedano ad Agent~3:

\begin{itemize}
  \item \textbf{Stato Binario Certificato:} Il sistema distingue nettamente
        tra \texttt{VALID} (0 issue, procede ad Agent~3) e \texttt{ISSUES\_FOUND}
        (problemi rilevati, loop continua). Nessun modello parzialmente validato
        può passare al nodo successivo.

  \item \textbf{Segregazione delle Funzioni:} La \textit{detection} (Step~1--2)
        è separata dalla \textit{resolution} (Step~0 + Step~4), garantendo
        che le modifiche siano sempre motivate da issue specifici rilevati in
        modo indipendente.

  \item \textbf{Tracciabilità delle Decisioni:} I metadati
        \texttt{\_userGuidedFixes} e \texttt{\_ownershipDecisions} accumulano
        la storia di tutte le scelte del business stakeholder durante il loop,
        fornendo un audit trail delle motivazioni architetturali.
\end{itemize}

\section{Tracciabilità e Audit}

\begin{table}[H]
\centering
\renewcommand{\arraystretch}{1.3}
\begin{tabularx}{\textwidth}{l l X}
\toprule
\textbf{Strumento} & \textbf{Scopo} & \textbf{Dettaglio Operativo} \\
\midrule
Task ID (UUID) & Tracciabilità E2E & Collega ogni log all'analisi originante. \\
\texttt{\_userGuidedFixes} & Audit delle decisioni & Storia delle risoluzioni manuali nel modello. \\
\texttt{output\_agent/} & Persistenza risultati & JSON completo di ogni analisi per Agent~3. \\
Log strutturati & Diagnostica & Logging JSON con level, timestamp e task\_id. \\
\bottomrule
\end{tabularx}
\caption{Strumenti di Audit e Tracciabilità}
\label{tab:audit-tools}
\end{table}

% =============================================================================
\chapter{Osservabilità e Monitoraggio}
% =============================================================================

\section{Logging Strutturato}

Il backend produce log in formato strutturato.
Ogni riga di log include task ID,
livello, timestamp e messaggio descrittivo.

\begin{technicalbox}{Dettaglio Tecnico: Tracciabilità per Task}
Ogni step della pipeline è tracciato con il task\_id:
\begin{itemize}
  \item \texttt{[task\_id] Starting analysis pipeline...}
  \item \texttt{[task\_id] Found N semantic issues}
  \item \texttt{[task\_id] Found N conflict issues}
  \item \texttt{[task\_id] Generated N questions}
  \item \texttt{[task\_id] Suggested max iterations: N}
\end{itemize}
I log sono persistiti nel volume \texttt{agent2\_logs} e accessibili
tramite \texttt{docker logs agent2-api}.
\end{technicalbox}

\section{Health Check e Monitoring}

L'endpoint \texttt{GET /health} fornisce lo stato in tempo reale di tutti
i servizi interni. La risposta include:

\begin{itemize}
  \item \textbf{Status generale:} indicatore di salute del sistema (\texttt{healthy}/\texttt{unhealthy}).
  \item \textbf{Stato dei 4 servizi:} SemanticAnalyzer, ConflictDetector, QuestionGenerator e ModelRefiner, ciascuno con un flag di disponibilità.
  \item \textbf{Configurazione corrente:} URL di Ollama, modello LLM in uso e modello di embedding attivo.
\end{itemize}

Docker Compose esegue automaticamente questo health check ogni 30 secondi
(\texttt{interval: 30s}, \texttt{retries: 3}), garantendo il riavvio
automatico del container in caso di anomalia.
\section{Indicatori di Performance (KPI)}

\begin{table}[H]
\centering
\renewcommand{\arraystretch}{1.3}
\begin{tabularx}{\textwidth}{l l X}
\toprule
\textbf{KPI} & \textbf{Origine} & \textbf{Diagnostica} \\
\midrule
Latenza /analyze con LLM    & Log applicativi & Rilevamento colli di bottiglia Ollama. \\
Latenza /analyze senza LLM  & Log applicativi & Baseline performance regole DDD. \\
Iterazioni per convergenza  & ValidationResponse & Misura efficacia del raffinamento. \\
LLM Fallback rate           & Log ``using heuristic'' & Monitoraggio disponibilità Ollama. \\
Issue per categoria         & summary JSON & Distribuzione tipologie di errori DDD. \\
\bottomrule
\end{tabularx}
\caption{Indicatori di Performance e Diagnostica}
\label{tab:kpi}
\end{table}

\section{Monitoraggio Kafka (Predisposto)}

Con l'attivazione del \texttt{docker-compose.kafka.yml}, Kafka UI
(porta 8080) fornisce una dashboard web per:
\begin{itemize}
  \item Ispezione dei topic \texttt{agent1-output} e \texttt{agent2-output}
  \item Visualizzazione dei messaggi in transito
  \item Monitoraggio del consumer group \texttt{agent2-consistency-analyzer}
  \item Rilevamento del consumer lag per diagnosi di ritardi nella pipeline
\end{itemize}

% =============================================================================
\chapter{Deployment e Operatività}
% =============================================================================

\section{Topologia di Esecuzione}

Il sistema è distribuito come uno stack di container orchestrati tramite
Docker Compose (cfr.\ Tabelle~\ref{tab:core-services} e~\ref{tab:kafka-services}
nel Capitolo~6 per il dettaglio dei singoli servizi), disponibile in due
configurazioni principali:

\begin{table}[H]
\centering
\small
\renewcommand{\arraystretch}{1.35}
\setlength{\tabcolsep}{10pt}

\begin{tabularx}{\textwidth}{
  >{\raggedright\arraybackslash}p{6cm}
  @{\hspace{5pt}}
  >{\raggedright\arraybackslash}p{4cm}
  @{\hspace{5pt}}
  >{\raggedright\arraybackslash}X
}
\toprule
\textbf{File} & \textbf{Container} & \textbf{Quando usarlo} \\
\midrule
\texttt{docker-compose.yml}
& 2 (agent2-api, n8n)
& Configurazione standard per il flusso interattivo REST \\

\texttt{docker-compose.kafka.yml}
& +4 (zookeeper, kafka, consumer, kafka-ui)
& Da usare solo per pipeline automatizzate agent-to-agent (A2A) \\
\bottomrule
\end{tabularx}

\caption{Configurazioni Docker Compose del sistema}
\label{tab:docker-configs}
\end{table}
\begin{technicalbox}{Dettaglio tecnico: avvio e healthcheck}
\begin{lstlisting}[style=bashbox]
# flusso standard (2 container, ~3 GB RAM)
docker compose up -d --build

# flusso completo con kafka (6 container, ~9 GB RAM)
docker compose -f docker-compose.yml \
  -f docker-compose.kafka.yml up -d --build

# verifica stato servizio
curl http://localhost:8002/health
\end{lstlisting}

\vspace{2pt}
Il Dockerfile utilizza una build multi-stage per ridurre la dimensione finale
dell'immagine: il primo stage compila le dipendenze, mentre il secondo include
esclusivamente l'ambiente di runtime e l'applicazione.
\end{technicalbox}

\section{Configurazione e Variabili d'Ambiente}

Il comportamento di Agent~2 è interamente configurabile tramite variabili
d'ambiente, con default hardcoded nel \texttt{docker-compose.yml}:

\begin{lstlisting}[language=bash]
# Ollama
OLLAMA_BASE_URL=http://host.docker.internal:11434
OLLAMA_MODEL=qwen2.5:14b
OLLAMA_TIMEOUT=180

# Analisi
EMBEDDINGS_MODEL=all-MiniLM-L6-v2
ENTITY_OVERLAP_THRESHOLD=0.85
SEMANTIC_SIMILARITY_THRESHOLD=0.75
MAX_FOLLOW_UP_QUESTIONS=5

# Kafka (futuro)
KAFKA_BOOTSTRAP_SERVERS=kafka:9092
KAFKA_INPUT_TOPIC=agent1-output
KAFKA_OUTPUT_TOPIC=agent2-output
KAFKA_GROUP_ID=agent2-consistency-analyzer

# Logging
LOG_LEVEL=INFO
\end{lstlisting}

\section{Persistenza e Gestione dei Dati}

\begin{itemize}
  \item \textbf{\texttt{./knowledge\_base}} $\rightarrow$ \texttt{/app/knowledge\_base:ro}:
        Regole DDD e checklist di validazione (sola lettura).
  \item \textbf{\texttt{./input\_agent}} $\rightarrow$ \texttt{/app/input\_agent:ro}:
        File JSON di esempio per test senza Agent~1 (sola lettura).
  \item \textbf{\texttt{agent2\_output}} (volume Docker): Output delle analisi
        in \texttt{/app/output\_agent}, consumato da Agent~3.
  \item \textbf{\texttt{agent2\_logs}} (volume Docker): Log dell'applicazione
        in \texttt{/app/logs}.
\end{itemize}
\section{Checklist di Manutenzione Operativa}

\begin{table}[H]
\centering
\renewcommand{\arraystretch}{1.3}
\begin{tabularx}{\textwidth}{l l X}
\toprule
\textbf{Attività} & \textbf{Strumento} & \textbf{Obiettivo} \\
\midrule
Health Check        & \texttt{GET /health}         & Verifica disponibilità backend e servizi. \\
Stato container     & \texttt{docker ps}           & Verifica che tutti i container siano healthy. \\
Log analisi         & \texttt{docker logs agent2-api} & Diagnostica errori LLM o pipeline. \\
Output archiviati   & \texttt{output\_agent/}      & Verifica persistenza dei modelli validati. \\
Modello Ollama      & \texttt{ollama list}         & Verifica che \texttt{qwen2.5:14b} sia disponibile. \\
\bottomrule
\end{tabularx}
\caption{Checklist Operativa di Manutenzione}
\label{tab:maintenance}
\end{table}

\section{Garanzia della Qualità e Testing}

\begin{technicalbox}{Dettaglio Tecnico: File di Test Predisposti}
La directory \texttt{input\_agent/} contiene 5 file JSON con copertura
sistematica di tutte le tipologie di issue:

\begin{center}
\begin{tabularx}{\textwidth}{l l X}
\toprule
\textbf{File} & \textbf{Dominio} & \textbf{Issue / Scopo} \\
\midrule
\texttt{example\_good.json}    & E-Commerce & Nessuno -- verifica \texttt{VALID} \\
\texttt{example\_demo.json}    & E-Commerce & 3 bug intenzionali -- demo rapida \\
\texttt{example\_bad.json}     & E-Commerce & Stress test completo (tutti i tipi di issue) \\
\texttt{case\_healthcare.json} & Sanità     & 4 issue clinici (overlap Patient, ambiguity Visita) \\
\texttt{case\_banking.json}    & Banking    & 4 issue finanziari (naming, pattern, conflict) \\
\bottomrule
\end{tabularx}
\end{center}
Gli endpoint \texttt{/source/files}, \texttt{/source/file} e
\texttt{/source/analyze-by-file} permettono di testare direttamente
questi file senza Agent~1.
\end{technicalbox}

\section{Resilienza Operativa}

Il deployment è progettato per essere self-contained su singola macchina.
L'esternalizzazione di Ollama sull'host permette di:
\begin{itemize}
  \item Aggiornare il modello LLM (\texttt{ollama pull}) senza riavviare
        i container
  \item Sperimentare con modelli diversi modificando solo
        \texttt{OLLAMA\_MODEL} nel \texttt{.env}
  \item Scalare il consumer Kafka modificando \texttt{CONSUMER\_REPLICAS}
        senza rebuild dell'immagine
\end{itemize}

\section{Backup e Disaster Recovery}

\begin{itemize}
  \item \textbf{Knowledge Base:} I file \texttt{ddd\_rules.md} e
        \texttt{validation\_checklist.json} sono versionati in Git.
        Nessun backup aggiuntivo richiesto.

  \item \textbf{Output Analisi:} Il volume \texttt{agent2\_output} contiene
        i risultati persistiti. In caso di guasto, è sufficiente riavviare
        lo stack: il volume Docker sopravvive ai container.

  \item \textbf{Workflow n8n:} I file \texttt{n8n/*.json} sono esportazioni
        versionabili dei workflow. Il volume \texttt{n8n\_data} conserva
        la configurazione n8n tra i riavvii.

  \item \textbf{Kafka (futuro):} Con retention configurata a 168 ore (7 giorni),
        i messaggi non ancora consumati da Agent~3 sono recuperabili via
        offset reset anche dopo un riavvio del consumer.
\end{itemize}

% =============================================================================
\chapter{Limiti e Scelte Progettuali}
% =============================================================================

\section{Scelte di Semplificazione (Trade-offs)}

Per mantenere Agent~2 autonomo, eseguibile su una singola macchina e
facilmente testabile, sono state adottate le seguenti scelte:

\begin{itemize}
  \item \textbf{Task Store In-Memory:} I task A2A sono conservati in una
        struttura dati in memoria. Non persistono tra i riavvii del container.
        Per produzione è prevista la sostituzione con Redis.

  \item \textbf{Ollama fuori da Docker:} Scelta deliberata per garantire
        accesso diretto alla GPU senza overhead di virtualizzazione.
        Introduce una dipendenza sull'host che deve essere gestita manualmente.

  \item \textbf{n8n come Frontend:} L'uso di n8n come orchestratore workflow
        evita lo sviluppo di un frontend React/Vue dedicato, riducendo la
        complessità dello stack. Il compromesso è la dipendenza dall'ecosistema n8n
        e la limitazione nell'accesso alle variabili d'ambiente in n8n v2.

  \item \textbf{Kafka Separato in File Distinto:} La separazione in
        \texttt{docker-compose.kafka.yml} risparmia ~6 GB RAM e 6 CPU nei
        casi d'uso standard (solo REST), al costo di un comando più complesso
        per l'avvio con la pipeline completa.
\end{itemize}

\begin{technicalbox}{Technical Insight: Filosofia ``Zero Dependency''}
Agent~2 funziona con un singolo comando (\texttt{docker compose up -d}),
senza richiedere database esterni, servizi cloud o identity provider.
L'unico prerequisito esterno è Ollama sull'host, giustificato dalla necessità
di accesso diretto alla GPU per prestazioni di inferenza accettabili.
\end{technicalbox}
\section{Analisi dei Compromessi Progettuali}

\begin{table}[H]
\centering
\renewcommand{\arraystretch}{1.3}
\begin{tabularx}{\textwidth}{l l X}
\toprule
\textbf{Area} & \textbf{Compromesso Accettato} & \textbf{Motivazione Tecnica} \\
\midrule
Pipeline LLM   & LLM-first con fallback euristico & Qualità massima senza degradare la disponibilità. \\
Iterazioni     & Soft cap adattivo (non hard stop)  & Convergenza garantita verso \texttt{VALID}. \\
AI locale      & Ollama sull'host (non in Docker)   & Accesso diretto GPU, privacy dei dati, autonomia. \\
Stato task     & In-memory (non persistente)        & Semplicità; sufficiente per sessioni di analisi. \\
Kafka          & Pipeline predisposta, non attiva   & Riduzione overhead nella modalità interattiva. \\
\bottomrule
\end{tabularx}
\caption{Matrice dei Compromessi Progettuali}
\label{tab:tradeoffs}
\end{table}

\section{Consapevolezza del Perimetro Progettuale}

Le scelte descritte definiscono Agent~2 come un sistema modulare e indipendente.
I limiti identificati non rappresentano debolezze strutturali, ma decisioni
deliberate per mantenere il software:

\begin{description}
  \item[Eseguibile:] Il sistema parte su una singola macchina con un solo comando.
  \item[Testabile:] I file di esempio in \texttt{input\_agent/} coprono tutti e 7
        i tipi di issue senza necessità di Agent~1.
  \item[Estendibile:] L'aggiunta di nuove regole DDD richiede solo la modifica di
        \texttt{ddd\_rules.md} e \texttt{validation\_checklist.json}, senza
        cambiamenti al codice.
\end{description}

% =============================================================================
\chapter{Roadmap ed Evoluzioni Future}
% =============================================================================

\section{Pipeline Kafka Automatica (Priorità Alta)}

La prima evoluzione naturale -- già predisposta architetturalmente -- è l'attivazione
della pipeline automatica Agent~1 $\rightarrow$ Agent~2 $\rightarrow$ Agent~3
via Kafka.

Il file \texttt{docker-compose.kafka.yml} è già completo e testato. L'attivazione
richiede:
\begin{enumerate}
  \item Avvio con il comando multi-file
  \item Configurazione di Agent~1 per pubblicare su \texttt{agent1-output}
  \item Configurazione di Agent~3 per consumare \texttt{agent2-output}
\end{enumerate}

\begin{technicalbox}{Technical Insight: Scalabilità del Consumer}
La variabile \texttt{CONSUMER\_REPLICAS} permette di scalare orizzontalmente
il consumer da 1 a $N$ istanze. Con 4 partizioni per topic e
\texttt{MAX\_WORKERS=4} per replica, il sistema può processare fino a 16 modelli
in parallelo, sufficiente per pipeline enterprise di medie dimensioni.
\end{technicalbox}

\section{Persistenza Task con Redis}

La sostituzione del task store in-memory con Redis permetterebbe:
\begin{itemize}
  \item Sopravvivenza dei task ai riavvii del container
  \item Condivisione dello stato tra più istanze dell'API in modalità scaled
  \item TTL automatico per pulizia dei task obsoleti
\end{itemize}
L'integrazione è predisposta architetturalmente: sostituire il dizionario
lo store in-memory con un client Redis richiede modifiche minime all'endpoint layer.
\section{Evoluzione della Knowledge Base}

La knowledge base DDD è progettata per crescere in modo indipendente dal codice:
\begin{itemize}
  \item \textbf{Regole Domain-Specific:} Aggiunta di regole specializzate per
        domini verticali (banking, healthcare, e-commerce) con pattern e
        anti-pattern settoriali.
  \item \textbf{Versionamento Semantico:} Introduzione di una versione per la
        \texttt{validation\_checklist.json} per tracciare l'evoluzione delle
        regole nel tempo.
  \item \textbf{Motore Esterno (Groq):} Per pipeline con requisiti di latenza
        ultra-bassa, integrazione di un Adapter AI capace di commutare tra
        Ollama (locale, privacy) e Groq (cloud, latenza $<$1s) in base al carico
        o alla complessità della richiesta.
\end{itemize}

\section{Integrazione con Agent 1 e Agent 3}

\begin{itemize}
  \item \textbf{Feedback verso Agent~1:} In caso di modello irrisolvibile
        (troppi issue strutturali dopo $N$ iterazioni), Agent~2 potrebbe
        rimandare il modello ad Agent~1 con annotazioni specifiche sui gap
        informativi da colmare nell'intervista.

  \item \textbf{Pre-processing per Agent~3:} Arricchimento del modello raffinato
        con metadati di specifica (hint per pattern architetturali, suggerimenti
        di tecnologie) per facilitare la generazione di microservizi da parte
        di Agent~3.
\end{itemize}

% =============================================================================
\chapter{Conclusioni}
% =============================================================================

\section{Sintesi dei Risultati}

La realizzazione di Agent~2 ha permesso di implementare un quality gate
intelligente per la pipeline \textit{Intelligent Domain Architect},
capace di trasformare un domain model grezzo in un artefatto certificato
e pronto per la generazione di microservizi.

Il progetto consegna un microservizio che risponde a tre sfide fondamentali
del moderno software design:

\begin{itemize}
  \item \textbf{Validazione Semantica Profonda:} La combinazione di embedding
        vettoriali (\texttt{all-MiniLM-L6-v2}) e analisi LLM (Qwen2.5~14B)
        rileva tipologie di incoerenza che sfuggono agli strumenti di analisi
        sintattica convenzionale.

  \item \textbf{Raffinamento Collaborativo:} Il loop iterativo n8n trasforma
        la validazione da processo passivo a dialogo attivo con il business
        stakeholder, garantendo che ogni modifica al modello sia consapevole
        e motivata.

  \item \textbf{Architettura Evolutiva:} La separazione in due file
        docker-compose, la stessa immagine Docker per API e consumer Kafka,
        e il sistema di fallback AI garantiscono che il sistema sia pronto
        per evoluzioni verso la pipeline automatica senza stravolgere
        l'architettura esistente.
\end{itemize}

\section{Valore Progettuale e Coerenza}

La coerenza del lavoro risiede nel perfetto allineamento tra i requisiti
della pipeline e le scelte implementative. Ogni decisione tecnica -- dalla
scelta di Qwen2.5~14B rispetto a Llama3, al soft cap adattivo nel loop n8n,
alla separazione del Kafka in file distinto -- è guidata da un obiettivo
concreto e misurabile.

Agent~2 dimostra che la validazione automatica di modelli DDD richiede un
approccio ibrido: regole formali per i casi chiari, ragionamento LLM per
i casi complessi, e interazione umana per le decisioni architetturali
che richiedono contesto di business non codificabile.
\section{Considerazioni Finali}

\begin{tcolorbox}[enhanced, colback=lightgray, colframe=boxheader,
  title=\textbf{Sintesi del Valore Aggiunto}, fonttitle=\small\bfseries,
  colbacktitle=boxheader, coltitle=white]
Agent~2 non è un semplice linter DDD, ma un \textbf{mediatore intelligente}
tra la fase di analisi del dominio e quella di implementazione.
Il ciclo iterativo guidato dall'LLM assicura che il domain model
converga verso la coerenza senza imporre una visione rigida,
ma guidando il business stakeholder verso decisioni architetturalmente
solide attraverso domande contestuali e risposte suggerite azionabili.
\end{tcolorbox}

Il progetto conferma che un'architettura a microservizi orientata agli
agenti può implementare processi di governance complessi in modo modulare,
scalabile e trasparente, ponendo le basi per una pipeline di sviluppo
software sempre più automatizzata ma rigorosamente controllata.

% =============================================================================
\chapter*{Glossario}
\addcontentsline{toc}{chapter}{Glossario}
% =============================================================================

\begin{description}[style=nextline]
  \item[A2A Protocol (Agent-to-Agent):] Protocollo standardizzato (v0.3) per
        la comunicazione strutturata tra agenti AI, basato su Task, Message e Artifact.
  \item[ACL (Anticorruption Layer):] Pattern DDD in cui un bounded context
        costruisce uno strato di traduzione per proteggersi dal modello di un
        contesto esterno, evitando la contaminazione del proprio dominio.
  \item[Agent Card:] Documento JSON pubblicato su \texttt{/.well-known/agent.json}
        che descrive le capacità, gli skill e gli schemi di sicurezza di un agente A2A.
  \item[Aggregate:] Cluster di entity e value object trattato come unità
        transazionale unica nel DDD. L'accesso avviene tramite l'Aggregate Root,
        che garantisce le invarianti di business.
  \item[Bounded Context:] Confine logico all'interno di un dominio DDD in cui
        un modello e un linguaggio ubiquitario hanno un significato preciso e coerente.
  \item[Consumer Group:] In Kafka, gruppo di consumer che si ripartiscono le
        partizioni di un topic per parallelizzare la lettura. Se un consumer
        muore, le sue partizioni vengono riassegnate automaticamente.
  \item[Context Map:] Rappresentazione delle relazioni tra bounded context nel DDD,
        con pattern di integrazione quali Shared Kernel, Customer-Supplier, ACL,
        Conformist, Open Host Service e Published Language.
  \item[Coreografia:] Pattern di coordinamento event-driven in cui nessun
        controllore centrale dirige il flusso: ogni servizio reagisce
        autonomamente agli eventi pubblicati dagli altri.
  \item[CQRS (Command Query Responsibility Segregation):] Pattern che separa
        il modello di scrittura (command) dal modello di lettura (query),
        consentendo di ottimizzare e scalare ciascuno indipendentemente.
  \item[DDD (Domain-Driven Design):] Approccio alla progettazione software che
        allinea l'implementazione al modello di business tramite un linguaggio condiviso.
  \item[Domain Event:] Fatto significativo accaduto nel dominio, espresso al
        tempo passato (es.\ \texttt{OrderPlaced}). Meccanismo naturale di
        comunicazione disaccoppiata tra bounded context.
  \item[Embedding:] Rappresentazione vettoriale numerica di un testo in uno
        spazio semantico ad alta dimensionalità, usata per calcolare similarità.
  \item[Entity Overlap:] Violazione DDD in cui la stessa entità appare in più
        bounded context senza differenziazione di ruolo o ownership.
  \item[Event Sourcing:] Pattern in cui lo stato di un'entità viene memorizzato
        come sequenza ordinata di eventi immutabili, anziché come snapshot.
        Lo stato corrente si ottiene riproducendo la sequenza.
  \item[FastAPI:] Framework per API REST asincrone ad alte prestazioni,
        con generazione automatica della documentazione OpenAPI.
  \item[GQA (Grouped-Query Attention):] Ottimizzazione architetturale adottata da
        Qwen2.5~14B che riduce il costo computazionale dell'attenzione, aumentando
        la velocità di inferenza.
  \item[Kafka:] Piattaforma di event streaming distribuita basata su un log
        immutabile e append-only, organizzata in topic e partizioni, con
        replicazione Leader/Follower per tolleranza ai guasti.
  \item[LLM-first:] Approccio che usa il Large Language Model come prima scelta
        per task di ragionamento, con fallback a euristica o template in caso di errore.
  \item[Qwen2.5~14B:] Modello linguistico a 14,8 miliardi di parametri di Alibaba Cloud,
        ottimizzato per output strutturato JSON, ragionamento multi-step e seguire
        istruzioni complesse. Usato via Ollama in formato Q4\_K\_M ($\sim$9 GB).
  \item[n8n:] Workflow automation engine open-source con interfaccia visuale,
        usato come orchestratore del loop interattivo utente.
  \item[Ollama:] Runtime locale per modelli LLM open-source, con supporto GPU
        nativo e API REST compatibile OpenAI.
  \item[Orchestrazione:] Pattern di coordinamento in cui un componente centrale
        conosce e dirige l'intero flusso di lavoro, invocando i servizi nella
        sequenza corretta e gestendo errori e compensazioni.
  \item[Partizione:] Suddivisione di un topic Kafka in segmenti indipendenti
        distribuiti su più broker, che consentono parallelismo in lettura e
        scrittura. L'ordine è garantito solo all'interno di una partizione.
  \item[RAG (Retrieval-Augmented Generation):] Tecnica che inietta documenti
        rilevanti nel prompt LLM per ancorare la generazione a fonti verificate.
        Usata nel conflict detector con le regole DDD come contesto.
  \item[Ubiquitous Language:] Vocabolario condiviso tra sviluppatori e domain expert
        all'interno di un bounded context, principio fondamentale del DDD.
  \item[Value Object:] Oggetto DDD immutabile definito dai suoi attributi,
        privo di identità propria. Due value object con gli stessi attributi
        sono considerati equivalenti (es.\ un indirizzo, un importo monetario).
  \item[Zookeeper:] Servizio di coordinamento distribuito richiesto da Kafka
        per la gestione dei broker e dei metadati dei topic.
\end{description}

% =============================================================================
\appendix
\chapter{Guida Rapida all'Installazione}
% =============================================================================

\section{Prerequisiti}

\begin{enumerate}
  \item \textbf{Docker Desktop} installato e in esecuzione
  \item \textbf{Ollama} installato sull'host con modello Qwen2.5~14B:
\begin{lstlisting}[language=bash]
ollama pull qwen2.5:14b        # ~9 GB, richiede connessione internet
ollama serve   # avvia il server se non parte automaticamente
\end{lstlisting}
  \item \textbf{GPU NVIDIA} con almeno 8 GB VRAM (raccomandato: RTX~4070 / 8 GB)
        oppure CPU con $\geq$16 GB RAM per inferenza CPU (più lenta)
\end{enumerate}

\section{Avvio tramite Docker Compose}

\begin{lstlisting}[language=bash]
# Clona il repository
git clone <repository-url>
cd agent-consistency-analyzer

# Avvio flusso standard (REST + n8n)
docker compose up -d --build

# Attendi che i container siano healthy (~40s)
docker ps  # verifica STATUS = healthy

# Test rapido
curl http://localhost:8002/health
\end{lstlisting}

In alternativa, usa lo script PowerShell incluso che automatizza l'intero avvio
(verifica Docker, avvia Ollama se necessario, controlla il modello LLM,
build dell'immagine e avvio dei container):
\begin{lstlisting}[language={}]
powershell -ExecutionPolicy Bypass -File .\START_DEMO.ps1
\end{lstlisting}

\section{Avvio con Pipeline Kafka (Futuro)}

\begin{lstlisting}[language=bash]
docker compose \
  -f docker-compose.yml \
  -f docker-compose.kafka.yml \
  up -d --build
\end{lstlisting}

\section{Variabili d'Ambiente (file \texttt{.env})}

Creare un file \texttt{.env} nella root del progetto:

\begin{lstlisting}[language=bash]
OLLAMA_BASE_URL=http://host.docker.internal:11434
OLLAMA_MODEL=qwen2.5:14b
OLLAMA_TIMEOUT=180

EMBEDDINGS_MODEL=all-MiniLM-L6-v2
ENTITY_OVERLAP_THRESHOLD=0.85
SEMANTIC_SIMILARITY_THRESHOLD=0.75
MAX_FOLLOW_UP_QUESTIONS=5

LOG_LEVEL=INFO
API_PORT=8002
\end{lstlisting}

\section{Pagine Utili}

\begin{table}[H]
\centering
\renewcommand{\arraystretch}{1.3}
\begin{tabularx}{\textwidth}{l X}
\toprule
\textbf{URL} & \textbf{Contenuto} \\
\midrule
\texttt{http://localhost:8002/docs}                   & Swagger UI -- documentazione interattiva API \\
\texttt{http://localhost:8002/health}                 & Health check e stato servizi \\
\texttt{http://localhost:8002/.well-known/agent.json} & Agent Card (A2A Protocol) \\
\texttt{http://127.0.0.1:5678}                        & n8n -- workflow editor \\
\texttt{http://127.0.0.1:5678/webhook/agent2-start}   & Loop interattivo (workflow attivo) \\
\texttt{http://localhost:8080}                        & Kafka UI (solo con docker-compose.kafka.yml) \\
\bottomrule
\end{tabularx}
\caption{Riferimenti URL dell'ambiente di sviluppo}
\label{tab:urls}
\end{table}

% =============================================================================
\chapter{Documentazione API (OpenAPI)}
% =============================================================================

\section{Accesso alla Documentazione}

Grazie all'integrazione nativa FastAPI, la documentazione interattiva è
disponibile non appena lo stack Docker è avviato:

\begin{center}
\texttt{http://localhost:8002/docs}
\end{center}

\section{Riepilogo Endpoint}

\begin{table}[H]
\centering
\renewcommand{\arraystretch}{1.3}
\begin{tabular}{llp{7cm}}
\toprule
\textbf{Metodo} & \textbf{Path} & \textbf{Descrizione} \\
\midrule
GET  & \texttt{/}                              & Informazioni agente e lista endpoint \\
GET  & \texttt{/health}                        & Stato servizi e configurazione corrente \\
GET  & \texttt{/config}                        & Parametri di configurazione (debug) \\
POST & \texttt{/analyze}                       & Analisi completa del domain model \\
GET  & \texttt{/source/files}                  & Lista file JSON in \texttt{input\_agent/} \\
GET  & \texttt{/source/file?name=\ldots}       & Contenuto di un file di esempio \\
POST & \texttt{/source/analyze-by-file}        & Analisi selezionando un file per nome \\
GET  & \texttt{/.well-known/agent.json}        & Agent Card (A2A Protocol v0.3) \\
POST & \texttt{/a2a/message/send}              & Invia domain model via A2A \\
GET  & \texttt{/a2a/tasks}                     & Lista task A2A in memoria \\
GET  & \texttt{/a2a/tasks/\{task\_id\}}        & Stato di un task specifico \\
POST & \texttt{/a2a/tasks/\{task\_id\}/cancel} & Cancella un task in esecuzione \\
\bottomrule
\end{tabular}
\caption{Riepilogo degli endpoint REST e A2A}
\label{tab:endpoints}
\end{table}
\section{Funzionalità Swagger UI}

L'interfaccia Swagger (\texttt{/docs}) permette di:
\begin{itemize}
  \item \textbf{Esplorare:} Visualizzare tutti gli endpoint con schema
        di request/response.
  \item \textbf{Testare:} Eseguire chiamate reali direttamente dal browser
        con la funzione \textit{Try it out}, visualizzando codici HTTP e
        tempi di risposta.
  \item \textbf{Validare:} Verificare la struttura JSON del domain model
        tramite lo schema OpenAPI generato automaticamente.
  \item \textbf{Integrare:} Scaricare la specifica OpenAPI in formato JSON
        o YAML per generare client SDK in qualsiasi linguaggio.
\end{itemize}

La documentazione è sempre sincronizzata con i modelli dati del backend,
garantendo che nessun endpoint sia non documentato.

% =============================================================================
\chapter{Troubleshooting}
% =============================================================================

\section{Problemi di Avvio}

\subsection{Container agent2-api non diventa \texttt{healthy}}

\begin{technicalbox}{Diagnosi e Risoluzione}
\textbf{Sintomo:} \texttt{docker ps} mostra \texttt{unhealthy} o il container
si riavvia in loop.

\textbf{Causa più comune:} Ollama non è raggiungibile da dentro il container.

\textbf{Diagnosi:}
\begin{lstlisting}[language=bash]
docker logs agent2-api        # cerca errori di connessione
curl http://localhost:11434   # verifica che Ollama sia attivo sull'host
\end{lstlisting}

\textbf{Soluzioni:}
\begin{enumerate}
  \item Avvia Ollama: \texttt{ollama serve}
  \item Verifica che \texttt{OLLAMA\_BASE\_URL=http://host.docker.internal:11434}
        sia impostato nel \texttt{.env}
  \item Su Linux potrebbe essere necessario usare l'IP dell'host:
        \texttt{OLLAMA\_BASE\_URL=http://172.17.0.1:11434}
\end{enumerate}
\end{technicalbox}

\subsection{Errore \texttt{model 'qwen2.5:14b' not found}}

\begin{technicalbox}{Diagnosi e Risoluzione}
\textbf{Causa:} Il modello non è stato scaricato in Ollama.

\textbf{Soluzione:}
\begin{lstlisting}[language=bash]
ollama pull qwen2.5:14b        # ~9 GB, richiede connessione internet
ollama list                   # verifica che qwen2.5:14b sia presente
\end{lstlisting}
Per usare un modello alternativo, modifica \texttt{OLLAMA\_MODEL} nel
\texttt{.env} (es. \texttt{llama3:8b}, \texttt{phi3:mini}).
\end{technicalbox}

\subsection{n8n non carica il workflow}

\begin{technicalbox}{Diagnosi e Risoluzione}
\textbf{Sintomo:} \texttt{http://127.0.0.1:5678/webhook/agent2-start} restituisce 404.

\textbf{Causa:} Il workflow non è stato importato o attivato.

\textbf{Passi:}
\begin{enumerate}
  \item Apri \texttt{http://127.0.0.1:5678} nel browser
  \item Login con le credenziali di default (configurate nella prima esecuzione)
  \item Vai su \textit{Workflows} $\rightarrow$ \textit{Import from file}
  \item Carica \texttt{n8n/workflow\_complete\_loop.json}
  \item Attiva il workflow con il toggle in alto a destra
\end{enumerate}
\end{technicalbox}

\section{Problemi di Analisi}

\subsection{L'analisi restituisce sempre \texttt{VALID} (nessun issue)}

\begin{technicalbox}{Diagnosi e Risoluzione}
\textbf{Cause possibili:}
\begin{itemize}
  \item Il domain model inviato è genuinamente valido
  \item La soglia \texttt{ENTITY\_OVERLAP\_THRESHOLD} (default: 0.85) è troppo alta
  \item Il modello embeddings non è stato caricato correttamente
\end{itemize}

\textbf{Test di verifica:}
\begin{lstlisting}[language=bash]
# Testa con il file di esempio che contiene issue noti
curl -X POST http://localhost:8002/source/analyze-by-file \
     -H "Content-Type: application/json" \
     -d '{"filename": "example_bad.json", "use_llm": false}'
# Atteso: issue rilevati per ogni tipologia
\end{lstlisting}

\textbf{Abbassa le soglie nel .env se necessario:}
\begin{lstlisting}[language=bash]
ENTITY_OVERLAP_THRESHOLD=0.75
SEMANTIC_SIMILARITY_THRESHOLD=0.65
\end{lstlisting}
\end{technicalbox}
\subsection{Timeout LLM (\texttt{use\_llm: true} molto lento)}

\begin{technicalbox}{Diagnosi e Risoluzione}
\textbf{Sintomo:} L'endpoint \texttt{/analyze} risponde dopo 90--120 secondi
o va in timeout.

\textbf{Causa:} Ollama sta usando la CPU invece della GPU, o la GPU è satura.

\textbf{Diagnosi:}
\begin{lstlisting}[language=bash]
# Su Windows con GPU NVIDIA
nvidia-smi
# Verifica che ollama.exe usi la GPU (GPU-Util > 0 durante l'inferenza)
\end{lstlisting}

\textbf{Soluzioni:}
\begin{enumerate}
  \item Riavvia Ollama: assicurati che CUDA sia installato correttamente
  \item Usa \texttt{use\_llm: false} per analisi senza LLM (più veloce, meno accurato)
  \item Aumenta \texttt{OLLAMA\_TIMEOUT} nel \texttt{.env}:
        \texttt{OLLAMA\_TIMEOUT=300}
  \item Prova un modello più leggero: \texttt{OLLAMA\_MODEL=phi3:mini} (2.2 GB)
\end{enumerate}
\end{technicalbox}

\subsection{Domande generate in lingua inglese invece di italiano}

\begin{technicalbox}{Diagnosi e Risoluzione}
\textbf{Causa:} Il modello LLM risponde nella lingua del domain model ricevuto.
Se \texttt{Agent 1} produce il modello in inglese, le domande saranno in inglese.

\textbf{Soluzione:} Aggiungi esplicitamente l'istruzione di lingua nel prompt
del question generator modificando il campo \texttt{description} delle entità
nel domain model, oppure configura Agent~1 per produrre modelli con
\texttt{ubiquitousLanguage} in italiano.
\end{technicalbox}
\section{Problemi con Kafka (Pipeline Predisposta)}

\subsection{Consumer non riceve messaggi da \texttt{agent1-output}}

\begin{technicalbox}{Diagnosi e Risoluzione}
\textbf{Diagnosi:}
\begin{lstlisting}[language=bash]
# Verifica stato topic e consumer lag
docker exec kafka kafka-consumer-groups.sh \
  --bootstrap-server localhost:9092 \
  --describe --group agent2-consistency-analyzer

# Controlla i log del consumer
docker logs agent2-consumer
\end{lstlisting}

\textbf{Cause comuni:}
\begin{enumerate}
  \item Il topic \texttt{agent1-output} non esiste ancora: Agent~1 non ha ancora
        pubblicato nessun messaggio.
  \item Offset reset: se il consumer parte con \texttt{auto.offset.reset=earliest},
        rileggerà tutti i messaggi storici; con \texttt{latest} solo i nuovi.
  \item Kafka non è ancora pronto: attendi il health check
        (\texttt{kafka-ui} su porta 8080 mostra lo stato del cluster).
\end{enumerate}
\end{technicalbox}

\section{Diagnostica Rapida}

\begin{table}[H]
\centering
\small
\renewcommand{\arraystretch}{1.25}
\setlength{\tabcolsep}{11pt}

\begin{tabularx}{\textwidth}{
  >{\raggedright\arraybackslash}p{4cm}
  >{\raggedright\arraybackslash}p{5.2cm}
  >{\raggedright\arraybackslash}X
}
\toprule
\textbf{Sintomo} & \textbf{Comando diagnostico} & \textbf{Cosa cercare} \\
\midrule
Container non parte
& \texttt{docker logs agent2-api}
& errori di avvio dell'applicazione \\

Ollama non raggiungibile
& \texttt{curl http://localhost:11434}
& risposta che indica che il servizio è attivo \\

Analisi lenta
& \texttt{nvidia-smi}
& utilizzo GPU durante l'inferenza \\

n8n webhook 404
& \texttt{docker logs n8n}
& workflow non attivato o webhook non registrato \\

Output non salvato
& \texttt{docker exec agent2-api ls /app/output\_agent/}
& presenza dei file JSON attesi \\

Kafka lag alto
& Kafka UI (\texttt{http://localhost:8080})
& consumer group lag elevato \\
\bottomrule
\end{tabularx}

\caption{Guida rapida alla diagnostica}
\label{tab:troubleshooting}
\end{table}


% -----------------------------------------------------------------------------
\end{document}
